
\setcounter{chapter}{15}

\chapter{Artin 环、离散赋值环与 Dedekind 整环}

本章始终设 \(R\) 是满足 \(1 \neq 0\) 的交换环。

\section{Artin 环}

本节我们研究满足理想降链条件（D.C.C.）的交换环的基本理论，即 \textbf{Artin 环}（以 E. Artin 命名）。人们起初也许会期望这些环具有与满足升链条件（A.C.C.）的交换环（即诺特环）类似的性质，但事实并非如此。Artin 环的结构受到很强的限制；例如 Artin 环必然也是诺特环（定理 3）。非交换 Artin 环在表示论中起核心作用（参见第 18、19 章）。

\begin{definition}
交换环 \(R\) 的\textbf{Krull 维数}（或简称\textbf{维数}）是 \(R\) 中互不相同的素理想的链 \(P_0 \subsetneq P_1 \subsetneq P_2 \subsetneq \cdots \subsetneq P_n\) 的可能最大长度。若 \(R\) 有任意长的互不相同素理想链，则称 \(R\) 的维数是无限的。
\end{definition}

有限维的环必然同时满足素理想上的升链与降链条件（虽然不必满足所有理想上的这两个条件）。域的维数为 0，而一个不是域的主理想整环的维数为 1。

我们很快会看到，满足理想上 D.C.C. 的环 \(R\) 的维数总是 0（即素理想都是极大理想）。若 \(R\) 是域 \(k\) 上的有限生成 \(k\)-代数且是整环，则 \(R\) 的维数等于 \(R\) 的分式域在 \(k\) 上的超越次数（参见习题 11）。特别地，这个维数与前面为仿射簇的维数引入的定义一致。上面定义的优点是它不涉及任何 \(k\)-代数结构，并且适用于任意交换环 \(R\).

\begin{definition}
\(R\) 的\textbf{Jacobson 根}是 \(R\) 中所有极大理想的交，记作 \(\operatorname{Jac}R\).
\end{definition}

Jacobson 根类似于群的 Frattini 子群，并有相应的性质（参见第 6.1 节习题 24）：

\begin{proposition}\label{pro:16.1}
设 \(\mathcal{J}\) 是交换环 \(R\) 的 Jacobson 根。

（1）若 \(I\) 是 \(R\) 的真理想，则 \((I, \mathcal{J})\)（由 \(I\) 与 \(\mathcal{J}\) 生成的理想）也是真理想。

（2）Jacobson 根包含 \(R\) 的幂零根：\(\operatorname{rad}0 \subseteq \operatorname{Jac}R\).

（3）元素 \(x\) 属于 \(\mathcal{J}\) 当且仅当对所有 \(r \in R\)，\(1-rx\) 是单位。

（4）（\textbf{Nakayama 引理}）若 \(M\) 是任意有限生成 \(R\)-模且 \(\mathcal{J}M = M\)，则 \(M = 0\).
\end{proposition}

\begin{proof}
若 \(I\) 是 \(R\) 的真理想，则对某个极大理想 \(M\) 有 \(I \subseteq M\). 由于 \(\mathcal{J} \subseteq M\)，也有 \((I, \mathcal{J}) \subseteq M\)，这证明了（1）。

（2）由两种根的定义以及第 15.2 节命题 12 推出，因为极大理想都是素理想。

设 \(1-rx\) 不是单位，\(M\) 是包含 \(1-rx\) 的极大理想。由于 \(1 \notin M\)，有 \(rx \notin M\)，故 \(x\) 不可能属于 \(\mathcal{J}\)（因为 \(\mathcal{J} \subseteq M\)）。反之，设 \(x \notin \mathcal{J}\)，即存在极大理想 \(M\) 使 \(x \notin M\). 则 \(R = (x, M)\)，故对某个 \(y \in M\) 有 \(1 = rx+y\). 于是 \(1-rx = y \in M\)，故 \(1-rx\) 不是单位，这证明了（3）。

为证明（4），假设 \(M \neq 0\)，设 \(n\) 是使 \(M\) 可由 \(n\) 个元素生成的最小整数，比如由 \(m_1, \ldots, m_n\) 生成。由 \(M = \mathcal{J}M\) 得
\[
m_n = r_1m_1+\cdots+r_nm_n
\]
（某些 \(r_1, r_2, \ldots, r_n \in \mathcal{J}\)）。于是 \((1-r_n)m_n = r_1m_1+\cdots+r_{n-1}m_{n-1}\). 由（3），\(1-r_n\) 是单位，故 \(m_n\) 属于由 \(m_1, \ldots, m_{n-1}\) 生成的模，与 \(n\) 的最小性矛盾。故 \(M = 0\)，证明完毕。
\end{proof}

\begin{definition}
交换环 \(R\) 称为\textbf{Artin 环}，或满足理想上的\textbf{降链条件}（简记 D.C.C.），若 \(R\) 中不存在理想的无限递减链，即只要 \(I_1 \supseteq I_2 \supseteq I_3 \supseteq \cdots\) 是 \(R\) 的理想的递减链，就存在正整数 \(m\) 使对所有 \(k \geq m\) 有 \(I_k = I_m\). 类似地，\(R\)-模 \(M\) 称为 \textbf{Artin} 的，若它满足子模上的 D.C.C.
\end{definition}

由格同构定理立即得到：Artin 环 \(R\) 对理想 \(I\) 的每个商环 \(R/I\) 也是 Artin 环。

下面关于 Artin 环的结果与定理 15.2 中的结果平行，其证明完全类似，故留作习题。

\begin{proposition}\label{pro:16.2}
下述命题等价：

（1）\(R\) 是 Artin 环；

（2）\(R\) 的每个非空理想集合在包含关系下有极小元。
\end{proposition}

下一个结果给出 Artin 环的主要结构定理。

\begin{theorem}\label{thm:16.3}
设 \(R\) 是 Artin 环。

（1）\(R\) 中只有有限多个极大理想。

（2）商环 \(R/(\operatorname{Jac}R)\) 是有限多个域的直积。更精确地，若 \(M_1, \ldots, M_n\) 是 \(R\) 的（有限多个）极大理想，则
\[
R/(\operatorname{Jac}R) \cong k_1 \times \cdots \times k_n,
\]
其中 \(k_i\) 是域 \(R/M_i\)（\(1 \leq i \leq n\)）。

（3）\(R\) 的每个素理想都是极大理想，即 \(R\) 的 Krull 维数为 0。\(R\) 的 Jacobson 根等于 \(R\) 的幂零根，且是幂零理想：对某个 \(m \geq 1\) 有 \((\operatorname{Jac}R)^m = 0\).

（4）\(R\) 同构于有限多个 Artin 局部环的直积。

（5）每个 Artin 环都是诺特环。
\end{theorem}

\begin{proof}
为证明（1），设 \(\mathcal{S}\) 是 \(R\) 中所有为有限多个极大理想之交的理想构成的集合。由命题 2，\(\mathcal{S}\) 有极小元，设为 \(M_1 \cap M_2 \cap \cdots \cap M_n\). 则对任何极大理想 \(M\) 有
\[
M \cap M_1 \cap M_2 \cap \cdots \cap M_n = M \cap M_1 \cap M_2 \cap \cdots \cap M_n,
\]
故 \(M \supseteq M_1 \cap M_2 \cap \cdots \cap M_n\). 由第 7.4 节习题 11，对某个 \(i\) 有 \(M \supseteq M_i\). 于是 \(M = M_i\)，故 \(M_1, \ldots, M_n\) 就是 \(R\) 的全部极大理想。

（2）的证明由把中国剩余定理（第 7.6 节）应用于 \(M_1, \ldots, M_n\) 立即推出，因为这些极大理想显然两两互余且它们的交是 \(\operatorname{Jac}R\).

对（3），我们先证明 \(\mathcal{J} = \operatorname{Jac}R\) 是幂零的。由 D.C.C.，存在某个 \(m > 0\) 使对所有正整数 \(i\) 有 \(\mathcal{J}^m = \mathcal{J}^{m+i}\). 反设 \(\mathcal{J}^m \neq 0\). 设 \(\mathcal{S}\) 是所有满足 \(I\mathcal{J}^m \neq 0\) 的真理想 \(I\) 构成的集合，故 \(\mathcal{J} \in \mathcal{S}\). 设 \(I_0\) 是 \(\mathcal{S}\) 的极小元。存在 \(x \in I_0\) 使 \(x\mathcal{J}^m \neq 0\)，故由极小性必有 \(I_0 = (x)\). 但此时 \(((x)\mathcal{J})\mathcal{J}^m = x\mathcal{J}^{m+1} = x\mathcal{J}^m\)，故由 \((x)\) 的极小性得 \((x) = (x)\mathcal{J}\). 由上面的 Nakayama 引理，\((x) = 0\)，矛盾。这证明 \(\operatorname{Jac}R\) 幂零。

由于 \(\operatorname{Jac}R\) 幂零，特别地 \(\operatorname{Jac}R \subseteq \operatorname{rad}0\)，故由命题 1 的第二个论断这两个理想相等。

\(R\) 中每个素理想 \(P\) 都含 \(R\) 的幂零根，从而由已证部分含 \(\operatorname{Jac}R\). \(P\) 的像是商环 \(R/(\operatorname{Jac}R) = k_1 \times \cdots \times k_n\) 中的素理想。但在环的直积 \(R_1 \times \cdots \times R_h\)（每个 \(R_j\) 含 1）中，每个理想都具有形式 \(I_1 \times \cdots \times I_h\)，其中 \(I_j\) 是 \(R_j\) 的理想（参见第 7.6 节习题 3）。由此可知 \(k_1 \times \cdots \times k_n\) 中的素理想由在一个分量中为 0 的元素组成。特别地这样的素理想在 \(k_1 \times \cdots \times k_n\) 中也是极大理想，从而 \(P\) 是 \(R\) 中的极大理想，这完成了（3）的证明。

设 \(M_1, \ldots, M_n\) 是 \(R\) 的全部互不相同的极大理想，如（3）中那样 \((\operatorname{Jac}R)^m = 0\). 则
\[
0 = (\operatorname{Jac}R)^m = (M_1 \cap \cdots \cap M_n)^m \supseteq M_1^m \cap \cdots \cap M_n^m \supseteq M_1^mM_2^m\cdots M_n^m = 0,
\]
故 \(M_1^m \cap \cdots \cap M_n^m = 0\). 由中国剩余定理可知
\[
R \cong (R/M_1^m) \times (R/M_2^m) \times \cdots \times (R/M_n^m),
\]
且每个 \(R/M_i^m\) 是以 \(M_i/M_i^m\) 为唯一极大理想的 Artin 环，这证明了（4）。

为证明（5），由（4）只需证明 Artin 局部环是诺特环，故设 \(R\) 是 Artin 环、唯一极大理想为 \(M\). 此时 \(M = \operatorname{Jac}R\)，故对某个正整数 \(m\) 有 \(M^m = (\operatorname{Jac}R)^m = 0\). 于是 \(R \cong R/M^m\)，而此时容易验证（参见习题 8）\(R/M^m\) 是诺特环当且仅当它是 Artin 环。
\end{proof}

\begin{corollary}\label{cor:16.4}
环 \(R\) 是 Artin 环当且仅当 \(R\) 是诺特环且 Krull 维数为 0.
\end{corollary}

\begin{proof}
正向蕴涵已在定理 3 中证明。现在设 \(R\) 是诺特环且 Krull 维数为 0，即 \(R\) 的素理想都是极大理想。由于 \(R\) 是诺特环，由第 15.2 节推论 22（3），理想 \((0) = P_1\cdots P_n\) 是（不必互不相同的）素理想之积，而这些素理想由于 \(R\) 的维数为 0 都是极大理想。由中国剩余定理，\(R\) 同构于有限多个形如 \(R/M^m\)（\(M\) 是 \(R\) 中的极大理想）的诺特环的直积。如定理（5）的证明中那样，\(R/M^m\) 是 Artin 环，由此 \(R\) 是 Artin 环。
\end{proof}

\noindent\textbf{例}

（1）设 \(n > 1\) 是整数。由于环 \(R = \mathbb{Z}/n\mathbb{Z}\) 有限，它是 Artin 环。若 \(n = p_1^{a_1}p_2^{a_2}\cdots p_s^{a_s}\) 是 \(n\) 分解为互不相同的素幂的唯一分解，则
\[
\mathbb{Z}/n\mathbb{Z} \cong (\mathbb{Z}/p_1^{a_1}\mathbb{Z}) \times (\mathbb{Z}/p_2^{a_2}\mathbb{Z}) \times \cdots \times (\mathbb{Z}/p_s^{a_s}\mathbb{Z}).
\]
每个 \(\mathbb{Z}/p_i^{a_i}\mathbb{Z}\) 是以 \((p_i)/(p_i^{a_i})\) 为唯一极大理想的 Artin 局部环，故这就是定理 3（4）给出的 \(\mathbb{Z}/n\mathbb{Z}\) 的分解。\(R\) 的 Jacobson 根是由 \(p_1p_2\cdots p_s\)（即 \(n\) 的无平方部分）生成的理想，而 \(R/(\operatorname{Jac}R) \cong (\mathbb{Z}/p_1\mathbb{Z}) \times \cdots \times (\mathbb{Z}/p_s\mathbb{Z})\) 是域的直积。由 \(p_i\)（\(i = 1, \ldots, s\)）生成的理想是 \(R\) 的极大理想。

（2）对任何域 \(k\)，作为 \(k\) 上向量空间有限维的 \(k\)-代数 \(R\) 是 Artin 环，因为 \(R\) 中的理想特别地是 \(R\) 的 \(k\)-子空间，故 \(R\) 中任何理想链的长度都被 \(\dim_kR\) 界住。

（3）设 \(f\) 是 \(k[x]\) 中的非零多项式（\(k\) 是域）。则由上一例，商环 \(R = k[x]/(f(x))\) 是 Artin 环。\(R\) 分解为 Artin 局部环之直积由下式给出：
\[
k[x]/(f(x)) \cong k[x]/(f_1(x)^{a_1}) \times \cdots \times k[x]/(f_s(x)^{a_s}),
\]
其中 \(f(x) = f_1(x)^{a_1}\cdots f_s(x)^{a_s}\) 是 \(f(x)\) 在 \(k[x]\) 中分解为互不相同的不可约多项式之幂的分解（参见第 9.5 节命题 16）。\(R\) 的 Jacobson 根是由 \(f(x)\) 的无平方部分生成的理想，而 \(R\) 的极大理想是由不可约因子 \(f_i(x)\)（\(i = 1, \ldots, s\)）生成的理想，与例 1 类似。

\subsection*{习题}

设 \(R\) 是含 1 的交换环，\(\mathcal{J}\) 是它的 Jacobson 根。

\begin{enumerate}
\item 设 \(R\) 是 Artin 环，\(I\) 是 \(R\) 中的理想。证明 \(R/I\) 也是 Artin 环。

\item 证明每个含 1 的有限交换环都是 Artin 环。

\item 证明 Krull 维数为 0 的整环是域。

\item 证明 Artin 整环是域。

\item 设 \(I\) 是 \(R\) 中的幂零理想且对某个 \(R\)-模 \(M\) 有 \(M = IM\). 证明 \(M = 0\).

\item 设 \(0 \to M' \to M \to M'' \to 0\) 是 \(R\)-模的正合序列。证明 \(M\) 是 Artin \(R\)-模当且仅当 \(M'\) 与 \(M''\) 都是 Artin \(R\)-模。

\item 设 \(R = F\) 是域。证明 \(R\)-模 \(M\) 是 Artin 模当且仅当它是诺特模，当且仅当 \(M\) 是 \(F\) 上有限维向量空间。

\item 设 \(M\) 是环 \(R\) 的极大理想，且对某个 \(n \geq 1\) 有 \(M^n = 0\). 证明 \(R\) 是诺特环当且仅当 \(R\) 是 Artin 环。〔注意过滤 \(R \supseteq M \supseteq \cdots \supseteq M^{n-1} \supseteq M^n = 0\) 中各相继商 \(M^i/M^{i+1}\)（\(i = 0, \ldots, n-1\)）都是域 \(F = R/M\) 上的模。然后利用前两习题与第 15.1 节习题 6。〕

\item 设 \(M\) 是有限生成 \(R\)-模。证明若 \(x_1, \ldots, x_n\) 是 \(M\) 的元素且它们在 \(M/\mathcal{J}M\) 中的像生成 \(M/\mathcal{J}M\)，则它们生成 \(M\). 由此证明若 \(R\) 是诺特环且 \(a_1, \ldots, a_n\) 在 \(\mathcal{J}/\mathcal{J}^2\) 中的像生成 \(\mathcal{J}/\mathcal{J}^2\)，则 \(\mathcal{J} = (a_1, \ldots, a_n)\).〔设 \(N\) 是由 \(x_1, \ldots, x_n\) 生成的子模，并对模 \(A = M/N\) 应用 Nakayama 引理。〕

\item 设 \(R = \mathbb{Z}_{(2)}\) 是 \(\mathbb{Z}\) 在素理想 \((2)\) 处的局部化。证明 \(\operatorname{Jac}R = (2)\) 是由 2 生成的理想。若 \(M = \mathbb{Q}\)，证明 \(M/2M\) 是有限生成 \(R\)-模，但 \(M\) 不是有限生成 \(R\)-模。为什么不与上一习题矛盾？〔注意 Nakayama 引理中的假设。〕

\item 设 \(V\) 是域 \(k\) 上的仿射簇，\(R = k[V]\) 是其坐标环。用 \(d_t(R)\) 表示分式域 \(k(V)\) 在 \(k\) 上的超越次数，\(d_P(R)\) 表示用素理想链定义的 \(R\) 的 Krull 维数。本习题证明 \(d_t(R) = d_P(R)\). 由 Noether 正规化引理，存在 \(R\) 的多项式子环 \(R_1 = k[y_1, \ldots, y_m]\) 使 \(R\) 在 \(R_1\) 上整。

（a）证明 \(d_t(R_1) = d_t(R) = m\) 且 \(d_P(R_1) = d_P(R)\). 由此我们可以假设 \(R = R_1\).〔用上升与下降定理（参见第 15.3 节定理 26）证明第二个等式。〕

（b）当 \(R = R_1\) 时，通过给出显式的长度为 \(m\) 的素理想链证明 \(d_P(R) \geq d_t(R)\).

（c）当 \(R = R_1\) 时，证明 \(R\) 的任何非零素理想都含有元素 \(f\) 使 \(R_{(f)}\) 在 \(R\) 上超越且超越次数为 1. 用归纳法证明 \(d_P(R) \leq d_t(R)\)，并由此 \(d_P(R) = d_t(R)\).

\item 设 \(R\) 是以 \(M\) 为极大理想的诺特局部环。

（a）商 \(M/M^2\) 是域 \(R/M\) 上的模（即向量空间）。证明 \(d = \dim_{R/M}(M/M^2)\) 有限。

（b）证明 \(M\) 作为 \(R\) 中的理想可以由 \(d\) 个元素生成，且不能由更少的元素生成。〔利用习题 9。〕

（c）设 \(R = k[x_1, \ldots, x_n]_{(x_1, \ldots, x_n)}\) 是域 \(k\) 上多项式环 \(k[x_1, \ldots, x_n]\) 在极大理想 \((x_1, \ldots, x_n)\) 处的局部化，\(M\) 是 \(R\) 中的极大理想。证明 \(\dim_{R/M}(M/M^2) = n = \dim R\).〔参见上一习题。〕

\noindent 可以证明对任何以 \(M\) 为极大理想的诺特局部环 \(R\) 有 \(\dim_{R/M}(M/M^2) \geq \dim R\). 若 \(\dim_{R/M}(M/M^2) = \dim R\)，则称诺特局部环 \(R\) 是\textbf{正则局部环}。事实上正则局部环必是整环，且必整闭。

\item 若 \(R\) 是诺特环，证明 \(\operatorname{Spec}R\) 上的 Zariski 拓扑是离散的（即每个子集既 Zariski 开又 Zariski 闭）当且仅当 \(R\) 是 Artin 环。

\item 设 \(I\) 是多项式环 \(k[x_1, x_2, x_3, \ldots]\) 中的理想 \((x_1, x_2^2, x_3^3, \ldots)\)（\(k\) 是域），\(R\) 是商环 \(k[x_1, x_2, \ldots]/I\). 证明理想 \((x_1, x_2, x_3, \ldots)\) 在 \(R\) 中的像是 \(R\) 中唯一的素理想，但它不是有限生成的。由此证明 \(R\) 是 Krull 维数为 0 的局部环，但不是 Artin 环。
\end{enumerate}

\section{离散赋值环}

上一节我们证明了 Artin 环就是 Krull 维数为 0 的诺特环。现在我们考虑最简单的 1 维诺特环，即第 8.1 节中首次引入的离散赋值环：

\begin{definition}
（1）域 \(K\) 上的一个\textbf{离散赋值}是满足下述条件的函数 \(v : K^\times \to \mathbb{Z}\)：

（i）\(v\) 是满射；

（ii）对所有 \(x, y \in K^\times\) 有 \(v(xy) = v(x)+v(y)\)；

（iii）对所有满足 \(x+y \neq 0\) 的 \(x, y \in K^\times\) 有 \(v(x+y) \geq \min\{v(x), v(y)\}\).

子环 \(\{x \in K \mid v(x) \geq 0\} \cup \{0\}\) 称为 \(v\) 的\textbf{赋值环}。

（2）整环 \(R\) 称为\textbf{离散赋值环}（D.V.R.），若 \(R\) 是 \(R\) 的分式域上某个离散赋值 \(v\) 的赋值环。
\end{definition}

通常把 \(v\) 通过定义 \(v(0) = +\infty\) 延拓到整个 \(K\)，此时（ii）与（iii）对所有 \(a, b \in K\) 成立。

\noindent\textbf{例}

（1）\(\mathbb{Z}\) 在任何非零素理想 \((p)\) 处的局部化 \(\mathbb{Z}_{(p)}\) 是关于如下定义的 \(\mathbb{Q}\) 上的离散赋值 \(v_p\) 的 D.V.R.（参见第 7.1 节习题 27）。\(\mathbb{Q}^\times\) 中每个元素 \(a/b\) 都可以唯一地写成 \(p^n(a_1/b_1)\) 的形式，其中 \(n \in \mathbb{Z}\)、\(a_1/b_1 \in \mathbb{Q}^\times\) 且 \(a_1\) 与 \(b_1\) 都与 \(p\) 互素。定义
\[
v_p\left(\frac{a}{b}\right) = v_p\left(p^n\frac{a_1}{b_1}\right) = n.
\]
容易验证 D.V.R. 的各条公理满足。我们称 \(v_p\) 为 \(\mathbb{Q}\) 上的 \textbf{\(p\)-进赋值}。相应的赋值环是满足 \(n \geq 0\) 的有理数连同 0，即分母不被 \(p\) 整除的有理数 \(a/b\)，这正是 \(\mathbb{Z}_{(p)}\).

（2）对任何域 \(F\)，设 \(f\) 是 \(F[x]\) 中的不可约多项式。域 \(F(x)\) 中每个非零元素都可以唯一地写成 \(f^n(a/b)\) 的形式，其中 \(n \in \mathbb{Z}\)、\(a/b \in F(x)^\times\) 且 \(a\) 与 \(b\) 都与 \(f\) 互素。则
\[
v_f\left(f^n\frac{a}{b}\right) = n
\]
定义了 \(F(x)\) 上的一个赋值，相应的赋值环是 \(F[x]\) 在 \(f\) 处的局部化 \(F[x]_f\)，由 \(F(x)\) 中分母不被 \(f\) 整除的有理函数组成。当 \(f = x-a\) 是 \(F[x]\) 中的 1 次多项式时，这个赋值给出元素在 \(x = a\) 处的零点阶（若 \(n \geq 0\)）或极点阶（若 \(n < 0\)）。

（3）以域 \(F\) 为系数的形式 Laurent 级数环 \(F((x))\) 上有由
\[
v\left(\sum_{i \geq n}a_ix^i\right) = n
\]
定义的离散赋值（参见第 7.2 节习题 5）。相应的 D.V.R. 是以 \(F\) 为系数的 \(x\) 的幂级数环 \(F[[x]]\).

注意 \(v(1) = v(1)+v(1)\) 蕴涵 \(v(1) = 0\)，故每个离散赋值环 \(R\) 都是满足 \(1 \neq 0\) 的含单位元的环。由于按定义 \(R\) 是域的子环，\(R\) 特别地是整环。容易看出 D.V.R. 是欧几里得整环（参见第 8.1 节例 4），故特别地也是主理想整环与唯一分解整环。事实上 D.V.R. 的因子分解与理想结构非常简单，如下一命题所示。

\begin{proposition}\label{pro:16.5}
设 \(R\) 是关于赋值 \(v\) 的离散赋值环，\(t\) 是 \(R\) 中满足 \(v(t) = 1\) 的任意元素。则

（1）非零元素 \(u \in R\) 是单位当且仅当 \(v(u) = 0\).

（2）每个非零元素 \(r \in R\) 都可以写成 \(r = ut^n\) 的形式（某个单位 \(u \in R\)、某个 \(n \geq 0\)）。\(R\) 的分式域中每个非零元素 \(x\) 都可以写成 \(x = ut^n\) 的形式（某个单位 \(u \in R\)、某个 \(n \in \mathbb{Z}\)）。

（3）\(R\) 的每个非零理想都是形如 \((t^n)\)（\(n \geq 0\)）的主理想。特别地，\(R\) 是诺特环。
\end{proposition}

\begin{proof}
若 \(u\) 是单位，则对某个 \(v \in R\) 有 \(uv = 1\)，于是 \(v(u)+v(v) = v(uv) = 1\)，而 \(v(u) \geq 0\)、\(v(v) \geq 0\)，这表明 \(v(u) = 0\)。反之，若 \(u \neq 0\) 且 \(v(u) = 0\)，则 \(u^{-1} \in K\) 满足 \(v(u^{-1})+v(u) = v(1) = 0\)，故 \(v(u^{-1}) = 0\) 且 \(u^{-1} \in R\)，即 \(u\) 是单位。这证明了（1）。

对（2），注意若 \(v(x) = n\) 则 \(v(xt^{-n}) = 0\)，故由（1）\(xt^{-n} = u\) 是 \(R\) 中的单位。于是 \(x = ut^n\)，其中 \(x \in R\) 当且仅当 \(n = v(x) \geq 0\).

若 \(I\) 是 \(R\) 中的非零理想，设 \(r \in I\) 是使 \(v(r)\) 最小的元素。若 \(v(r) = n\)，则由（2）\(r\) 与 \(t^n\) 相差一个单位，故 \(t^n \in I\) 且 \((t^n) \subseteq I\). 若 \(a\) 是 \(I\) 中任意非零元素，则由 \(n\) 的取法 \(v(a) \geq n\). 于是 \(v(at^{-n}) \geq 0\)，故 \(at^{-n} \in R\)，这表明 \(a \in (t^n)\). 因此 \(I = (t^n)\)，这证明了（3）的第一个论断。由此显然 \(R\) 中理想的升链都是有限的，即 \(R\) 是诺特环，证明完毕。
\end{proof}

\begin{definition}
若 \(R\) 是关于赋值 \(v\) 的 D.V.R.，则 \(R\) 中满足 \(v(t) = 1\) 的元素 \(t\) 称为 \(R\) 的\textbf{一致化参数}（或\textbf{局部参数}）。
\end{definition}

\begin{corollary}\label{cor:16.6}
设 \(R\) 是离散赋值环。

（1）\(R\) 是整闭局部环，其唯一极大理想由赋值严格大于 0 的元素给出：\(M = \{r \in R \mid v(r) > 0\}\). \(R\) 中每个非零理想都具有形式 \(M^n\)（某个整数 \(n \geq 0\)）。

（2）\(R\) 的素理想只有 \(M\) 与 \(0\)，即 \(\operatorname{Spec}R = \{0, M\}\). 特别地，D.V.R. 的 Krull 维数为 1.
\end{corollary}

\begin{proof}
任何唯一分解整环都在其分式域中整闭（第 15.3 节例 3），故 \(R\) 整闭。其余论断由命题 5 中 \(R\) 的理想的描述立即推出。
\end{proof}

离散赋值环的定义用分式域上的赋值表述得非常明确，因此从 \(R\) 的纯粹「内部」代数性质判断一个给定的环 \(R\) 是否是 D.V.R. 似乎很困难。事实上命题 5 与推论 6 中的环论性质刻画了离散赋值环。下面的定理给出离散赋值环的若干不显式提及赋值的等价代数描述。

\begin{theorem}\label{thm:16.7}
环 \(R\) 的下述性质等价：

（1）\(R\) 是离散赋值环；

（2）\(R\) 是具有唯一非零极大理想 \(P\) 的主理想整环；

（3）\(R\) 是具有唯一（在相伴意义下）不可约元素 \(t\) 的唯一分解整环；

（4）\(R\) 是诺特整环，且是唯一极大理想非零且为原理想的局部环；

（5）\(R\) 是诺特、整闭的整环，且是 Krull 维数为 1 的局部环，即 \(R\) 有唯一的非零素理想：\(\operatorname{Spec}R = \{0, M\}\).
\end{theorem}

\begin{proof}
（1）蕴涵其他各条性质已在上文证明。

若（2）成立，则由唯一分解整环中不可约元素生成素理想（第 8.3 节命题 12），（3）立即推出。

若（3）成立，则 \(R\) 中每个非零元素都可以唯一地写成 \(ut^n\) 的形式（某个单位 \(u\)、某个 \(n \geq 0\)）。于是 \(R\) 的分式域中每个非零元素都可以唯一地写成 \(ut^n\) 的形式（某个单位 \(u\)、某个 \(n \in \mathbb{Z}\)）。此时容易验证映射 \(v(ut^n) = n\) 是 \(R\) 的分式域上的离散赋值，且 \(R\) 是这个赋值的赋值环，故（1）成立。

设（4）成立，\(M = (t)\) 是 \(R\) 的唯一极大理想，\(M_0 = \bigcap_{i \geq 1}M^i\). 则 \(M_0 = MM_0\)，而由于 \(R\) 是诺特环，\(M_0\) 有限生成。由假设 \(M = \operatorname{Jac}R\)，故由 Nakayama 引理 \(M_0 = 0\). 若 \(I\) 是 \(R\) 的任意非零真理想，则存在 \(n \geq 0\) 使 \(I \subseteq M^n\) 但 \(I \nsubseteq M^{n+1}\). 取 \(a \in I-M^{n+1}\)，写 \(a = t^nu\)（某个 \(u \in R\)）。则 \(u \notin M\)，故 \(u\) 是局部环 \(R\) 中的单位。于是对每个 \(a \in I-M^{n+1}\) 有 \((a) = (t^n) = M^n\). 这表明 \(I = (t^n)\)，故 \(R\) 的每个理想都是主理想，即（2）成立。

我们已经证明了（1）、（2）、（3）、（4）等价，且其中每一条都蕴涵（5）。为完成证明，我们证明（5）蕴涵（4），即证明（5）中的理想 \(M\) 是主理想。由于 \(0 \neq M = \operatorname{Jac}R\) 且 \(M\) 由 \(R\) 的诺特性有限生成，由 Nakayama 引理（命题 1（4））有 \(M \neq M^2\). 取 \(t \in M-M^2\). 我们证明 \(M = (t)\). 由第 15.2 节命题 12，假设 \(M\) 是 \(R\) 中唯一的非零素理想意味着 \(M = \operatorname{rad}(t)\)，而第 15.2 节命题 14 又蕴涵 \(M\) 的某个幂包含在 \((t)\) 中。反设 \((t) \neq M\)，即 \(M^n \subseteq (t)\) 但 \(M^{n-1} \nsubseteq (t)\)（某个 \(n \geq 2\)）。则存在 \(x \in M^{n-1}-(t)\) 使 \(xM \subseteq (t)\). 注意 \(t \neq 0\)，故 \(y = x/t\) 属于 \(R\) 的分式域。又 \(y \notin R\)，因为 \(x = ty \notin (t)\). 然而由 \(x\) 的取法有 \(yM \subseteq R\)，且容易验证 \(yM\) 是 \(R\) 中的理想。若 \(yM = R\)，则对某个 \(m \in M\) 有 \(1 = ym\). 这将导致矛盾，因为此时 \(t = xm \in M^2\)，与 \(t\) 的取法矛盾。故 \(yM\) 是真理想，从而包含在 \(R\) 的唯一极大理想中，即 \(yM \subseteq M\). 现在 \(M\) 是有限生成 \(R\)-模，而 \(y\) 通过左乘作为 \(R\)-模同态作用在 \(M\) 上。与第 15.3 节命题 23 的证明中相同的（行列式）方法表明存在系数在 \(R\) 中的首一多项式 \(p\) 使对所有 \(m \in M\) 有 \(p(y)m = 0\). 由于 \(p(y)\) 是含 \(R\) 与 \(M\) 的域中的元素，必有 \(p(y) = 0\). 故 \(y\) 在 \(R\) 上整。由于假设 \(R\) 整闭，可知 \(y \in R\)，矛盾。故 \(M = (t)\) 是主理想，即（5）蕴涵（4），定理证明完毕。
\end{proof}

\begin{corollary}\label{cor:16.8}
若 \(R\) 是任何诺特、整闭的整环，\(P\) 是 \(R\) 的极小非零素理想，则 \(R\) 在 \(P\) 处的局部化 \(R_P\) 是离散赋值环。
\end{corollary}

\begin{proof}
由第 15.4 节的结果，局部化 \(R_P\) 是诺特环（命题 38（4））、整闭环（命题 49）、整环（命题 46（2）），并且是以唯一非零素理想为唯一素理想的局部环（命题 46（4）），故 \(R_P\) 满足定理中的（5）。
\end{proof}

\noindent\textbf{例}

（1）若 \(R\) 是任何主理想整环，则 \(R\) 在任一非零素理想 \(P = (p)\) 处的局部化 \(R_P\) 都是离散赋值环。这由推论 8 立即推出，因为 \(R\) 整闭（它是唯一分解整环，参见第 15.3 节例 3）且主理想整环中的非零素理想都是极大理想（第 8.2 节命题 7）。注意 \(R_P\) 的分式域 \(K\) 与 \(R\) 的分式域相同，故 \(R\) 中每个非零素元素 \(p\) 都给出 \(K\) 上的赋值 \(v_p\)，由公式
\[
v_p\left(p^n\frac{a}{b}\right) = n
\]
给出，其中 \(a\) 与 \(b\) 是不被 \(p\) 整除的 \(R\) 的元素。这推广了上面的例 1 与例 2。

（2）\(p\)-进整数环 \(\mathbb{Z}_p\) 是离散赋值环，因为它是具有唯一极大理想 \(p\mathbb{Z}_p\) 的主理想整环（参见第 7.6 节习题 11）。\(\mathbb{Z}_p\) 的分式域称为 \textbf{\(p\)-进数域}，记作 \(\mathbb{Q}_p\). 元素 \(p\) 是 \(\mathbb{Z}_p\) 的一致化参数，故 \(\mathbb{Q}_p\) 中每个非零元素都可以唯一地写成 \(p^nu\) 的形式（某个 \(n \in \mathbb{Z}\)、某个单位 \(u \in \mathbb{Z}_p^\times\)，其中 \(u = a_0+a_1p+a_2p^2+\cdots\) 且 \(0 < a_0 < p\)，如第 7.6 节习题 11（c）中那样）。相应的 \(\mathbb{Q}_p\) 上的 \(p\)-进赋值 \(v_p\) 于是由 \(v_p(p^nu) = n\) 给出。

域 \(K\) 上的离散赋值 \(v\) 如下定义一个相应的度量（或「距离函数」）\(d_v\)：固定任意正实数 \(\beta\)（\(\beta\) 的实际值对验证度量公理无关），对所有 \(a, b \in K\) 定义
\[
d_v(a, b) = \|a-b\|_v, \qquad \text{其中 } \|a\|_v = \beta^{-v(a)},
\]
并规定 \(d_v(a,a) = 0\). 容易验证 \(d_v\) 满足度量的三条公理：

（i）\(d_v(a,b) \geq 0\)，且等号成立当且仅当 \(a = b\)；

（ii）\(d_v(a,b) = d_v(b,a)\)，即 \(d_v\) 对称；

（iii）对所有 \(a, b, c \in K\) 有 \(d_v(a,b) \leq d_v(a,c)+d_v(c,b)\)，即 \(d_v\) 满足「三角不等式」。

三角不等式是离散赋值公理（iii）的推论。事实上更强的三角不等式成立：

（iii）′ 对所有 \(a, b, c \in K\) 有 \(d_v(a,b) \leq \max\{d_v(a,c), d_v(c,b)\}\).

因此 \(d_v\) 有时称为\textbf{超度量}。现在可以用 Cauchy 序列构造 \(K\) 关于 \(d_v\) 的完备化，记作 \(K_v\)，方式与由有理数 \(\mathbb{Q}\) 构造实数 \(\mathbb{R}\) 相同。不难证明 \(K_v\) 也是带有离散赋值的域，该赋值在 \(K\) 这个稠密子集上与 \(v\) 一致。

\noindent\textbf{例}

（1）考虑 \(\mathbb{Q}\) 上的 \(p\)-进赋值 \(v_p\)，取 \(\beta = p\). 记 \(\|a\|_v = \|a\|_p\)，故对与 \(p\) 互素的 \(a, b\) 有
\[
\left\|p^n\frac{a}{b}\right\|_p = p^{-n}.
\]
注意整数（或有理数）若被 \(p\) 的高次幂整除，则有小的 \(p\)-进绝对值。例如序列 \(1, p, p^2, p^3, \ldots\) 在 \(p\)-进度量下收敛到 0.

不难看出 \(\mathbb{Q}\) 关于 \(p\)-进度量的完备化是 \(p\)-进数域 \(\mathbb{Q}_p\)，而 \(\mathbb{Z}\) 的完备化是 \(p\)-进整数环 \(\mathbb{Z}_p\). 看出这一点的一种方法是验证完备化中的每个元素都可以表示为 \(p\)-进 Laurent 级数
\[
a = \sum_{n \geq n_0}a_np^n, \qquad \text{其中 } n_0 \geq 0,\ a_n \in \{0, 1, \ldots, p-1\},
\]
然后利用上面的例 2。用这个展开式表示，\(p\)-进赋值由 \(v_p(a) = n_0\)（当 \(a_{n_0} \neq 0\) 时）给出。

（2）类似地，\(F(x)\) 关于本节开头例 2 中赋值 \(v\) 的完备化给出同一组例子中例 3 里的域 \(F((x))\)，相应的赋值环是 \(F[[x]]\).

域 \(K\) 关于离散赋值的完备化是这样一个域 \(K_v\)，其中元素可以用一致化参数容易地描述。此外 \(K_v\) 是拓扑由度量 \(d_v\) 定义的拓扑空间。进一步，\(K_v\) 中元素的 Cauchy 序列收敛到 \(K_v\) 的元素（即 \(K_v\) 在 \(v\)-进拓扑下完备）。这与 \(\mathbb{Q}\) 关于通常 Euclid 度量的完备化 \(\mathbb{R}\) 的情形类似。这使得分析中的想法可以应用于这类环的研究，是研究代数数域与代数几何的重要工具。

\noindent\textbf{分式理想}

我们以用「分式理想」给出的离散赋值环的另一个刻画来结束本节。分式理想可以对任何整环定义：

\begin{definition}
对任何分式域为 \(K\) 的整环 \(R\)，\(R\) 的\textbf{分式理想}是 \(K\) 的 \(R\)-子模 \(A\)，使对某个非零 \(d \in R\) 有 \(dA \subseteq R\)（等价地，形如 \(d^{-1}I\) 的子模，其中 \(d \in R\) 非零、\(I\) 是 \(R\) 的理想）。
\end{definition}

这两个定义的等价性由 \(dA\) 是 \(R\) 的 \(R\)-子模（即理想）这一观察推出。

\(K\) 中分式理想的概念依赖于环 \(R\). 粗略地说，分式理想是相差一个固定的「分母」\(d\) 的 \(R\) 的理想。\(R\) 的理想也是 \(R\) 的分式理想（取分母 \(d = 1\)），并且是包含在 \(R\) 中的那些分式理想。为清楚起见，这些有时称为 \(R\) 的\textbf{整理想}。当 \(R\) 是诺特整环时，\(R\) 的分式理想与 \(K\) 的有限生成 \(R\)-子模相同（参见习题 6）。

对任何 \(x \in K\)，循环 \(R\)-模 \(Rx = \{rx \mid r \in R\}\) 称为由 \(x\) 生成的\textbf{主分式理想}。

若 \(A\) 与 \(B\) 是分式理想，它们的\textbf{乘积} \(AB\) 定义为形如 \(ab\)（\(a \in A\)、\(b \in B\)）的元素的所有有限和之集。若 \(A = d^{-1}I\)、\(B = (d')^{-1}J\)（\(I, J\) 是 \(R\) 的理想，\(d, d' \in R\) 非零），则 \(AB = (dd')^{-1}IJ\)，其中 \(IJ\) 是通常的乘积理想。特别地这表明两个分式理想的乘积是分式理想。

\begin{definition}
分式理想 \(A\) 称为\textbf{可逆}的，若存在分式理想 \(B\) 使 \(AB = R\)，此时称 \(B\) 是 \(A\) 的\textbf{逆}，记作 \(A^{-1}\).
\end{definition}

若 \(A\) 是可逆分式理想，则满足 \(AB = R\) 的分式理想 \(B\) 唯一：\(AB = AC = R\) 蕴涵 \(B = B(AC) = (BA)C = C\).

\begin{proposition}\label{pro:16.9}
设 \(R\) 是整环，\(A\) 是 \(R\) 的分式理想。

（1）若 \(A\) 是非零主分式理想，则 \(A\) 可逆。

（2）若 \(A\) 非零，则集合 \(A' = \{x \in K \mid xA \subseteq R\}\) 是 \(R\) 的分式理想。一般地 \(AA' \subseteq R\)，且 \(AA' = R\) 当且仅当 \(A\) 可逆，此时 \(A^{-1} = A'\).

（3）若 \(A\) 是 \(R\) 的可逆分式理想，则 \(A\) 有限生成。

（4）可逆分式理想之集在乘法下构成以 \(R\) 为单位元的阿贝尔群。非零主分式理想之集是这个可逆分式理想群的子群。
\end{proposition}

\begin{proof}
若 \(A = xR\) 是非零主分式理想，则取 \(B = x^{-1}R\) 表明 \(A\) 可逆，这证明了（1）。

容易看出 \(A'\) 是 \(K\) 的 \(R\)-子模。若 \(A\) 是非零分式理想，则存在非零元素 \(d \in R\) 使 \(dA \subseteq R\)，故 \(A\) 含有 \(R\) 的非零元素。设 \(a\) 是 \(A\) 中含于 \(R\) 的任意非零元素。则按 \(A'\) 的定义有 \(aA' \subseteq R\)，故 \(A'\) 是分式理想。同样按定义 \(AA' \subseteq R\). 若 \(AA' = R\) 则 \(A\) 可逆且逆为 \(A'\). 反之，若 \(AB = R\)，则由 \(A'\) 的定义有 \(B \subseteq A'\). 于是 \(R = AB \subseteq AA' \subseteq R\)，这表明 \(AA' = R\)，这证明了（2）。

若 \(A\) 可逆，则由（2）\(AA' = R\)，故对某些 \(a_1, \ldots, a_n \in A\) 与 \(a'_1, \ldots, a'_n \in A'\) 有 \(1 = a_1a'_1+\cdots+a_na'_n\). 若 \(a \in A\)，则 \(a = (aa'_1)a_1+\cdots+(aa'_n)a_n\)，其中每个 \(aa'_i \in R\)（由 \(A'\) 的定义）。由此 \(A\) 由 \(a_1, \ldots, a_n\) 在 \(R\) 上生成，故 \(A\) 有限生成，这证明了（3）。

最后，显然两个可逆分式理想的乘积仍可逆。这个乘积交换、结合，且对任何分式理想有 \(RA = A\). 可逆分式理想的逆按定义也是可逆分式理想，这证明了（4）的第一个论断。（4）的第二个论断立即推出，因为 \(xR\) 与 \(yR\) 的乘积是 \((xy)R\)，而 \(xR\) 的逆是 \(x^{-1}R\).
\end{proof}

\begin{definition}
若 \(R\) 是整环，则 \(R\) 的可逆分式理想群模去 \(R\) 的非零主分式理想子群所得的商群称为 \(R\) 的\textbf{类群}。\(R\) 的类群的阶称为 \(R\) 的\textbf{类数}。
\end{definition}

\(R\) 的类群是平凡群、类数为 1，当且仅当 \(R\) 是主理想整环。\(R\) 的类群衡量 \(R\) 的理想与主理想的接近程度。

\(R\) 的分式理想 \(A\) 是否可逆也与 \(A\) 作为 \(R\)-模是否投射有关。回忆 \(R\)-模 \(M\) 在 \(R\) 上投射当且仅当 \(M\) 是自由模的直和项（第 10.5 节命题 30）。等价地，\(M\) 投射当且仅当存在自由 \(R\)-模 \(F\) 与 \(R\)-模同态 \(f : F \to M\)、\(g : M \to F\) 使 \(f \circ g = 1\)（第 10.5 节命题 25）。

\begin{proposition}\label{pro:16.10}
设 \(R\) 是分式域为 \(K\) 的整环，\(A\) 是 \(R\) 的非零分式理想。则 \(A\) 可逆当且仅当 \(A\) 是投射 \(R\)-模。
\end{proposition}

\begin{proof}
先设 \(A\) 可逆，故如命题 9（2）中那样对某些 \(a_i \in A\)、\(a'_i \in A'\) 有 \(1 = \sum_{i=1}^{n}a_ia'_i\). 设 \(F\) 是以 \(y_1, \ldots, y_n\) 为基的自由 \(R\)-模。定义 \(f : F \to A\) 为 \(f(\sum_{i=1}^{n}r_iy_i) = \sum_{i=1}^{n}r_ia_i\)，\(g : A \to F\) 为 \(g(c) = \sum_{i=1}^{n}(ca'_i)y_i\). 显然 \(f\) 与 \(g\) 都是 \(R\)-模同态（注意由 \(A'\) 的定义 \(ca'_i \in R\)）。由于
\[
(f \circ g)(c) = f\left(\sum_{i=1}^{n}(ca'_i)y_i\right) = \sum_{i=1}^{n}(ca'_i)a_i = c\left(\sum_{i=1}^{n}a_ia'_i\right) = c,
\]
故 \(f \circ g = 1\)，即 \(A\) 是 \(F\) 的直和项，从而投射。

反之，设 \(A\) 非零且投射，故存在自由 \(R\)-模 \(F\) 与 \(R\)-同态 \(f : F \to A\)、\(g : A \to F\) 使 \(f \circ g = 1\). 固定 \(0 \neq a \in A\)，设 \(g(a) = \sum_{i=1}^{n}a_iy_i\)（\(a_i \in R\)，\(y_1, \ldots, y_n\) 是 \(F\) 的自由生成元的一部分）。定义 \(a_i = f(y_i)\) 与 \(a'_i = a_i/a \in K\)（\(i = 1, \ldots, n\)）。对任意 \(b \in A\)，由于 \(g\) 是 \(R\)-模同态有 \(bg(a) = ag(b) = g(ab)\). 把 \(g(b)\) 写成 \(g(b) = \sum_{i=1}^{n}b_iy_i+\sum_{j \in \mathcal{J}}b_jy_j\)，其中 \(\{y_j\}\)（\(j \in \mathcal{J}\)）是 \(F\) 的自由生成元集中的其余元素。则
\[
\sum_{i=1}^{n}ba_iy_i = g(ab) = \sum_{i=1}^{n}(ab)_iy_i+\sum_{j \in \mathcal{J}}(ab)_jy_j.
\]
我们可以在 \(F\) 的自由 \(R\)-模基上比较这个等式中元素的系数，由此 \(g(b) = \sum_{i=1}^{n}b_iy_i\)（\(b_i \in R\)）且 \(ba_i = ab_i\)（\(i = 1, \ldots, n\)）。特别地，由 \(a'_i\) 的定义可知对 \(A\) 的每个元素 \(b\)，\(ba'_i = b(a_i/a) = b_i\) 是 \(R\) 的元素。这表明 \(a'_i \in A'\)（\(i = 1, \ldots, n\)）。由于 \(f \circ g = 1\)，我们有
\[
a = a\left(\sum_{i=1}^{n}a_ia'_i\right) = \sum_{i=1}^{n}a_ia_i,
\]
于是 \(\sum_{i=1}^{n}a_ia'_i = 1\). 由此 \(AA' = R\)，故由命题 9 \(A\) 可逆，证明完毕。
\end{proof}

下一个结果表明，若整环 \(R\) 还是局部环，则分式理想是否可逆决定了 \(R\) 是否是 D.V.R.

\begin{proposition}\label{pro:16.11}
设整环 \(R\) 是局部环且不是域。则 \(R\) 是离散赋值环当且仅当 \(R\) 的每个非零分式理想都可逆。
\end{proposition}

\begin{proof}
若 \(R\) 是以 \(t\) 为一致化参数的 D.V.R.，则由命题 5，\(R\) 的每个非零理想都具有形式 \((t^n)\)（某个 \(n \geq 0\)），且 \(R\) 中每个元素 \(d\) 都可以写成 \(ut^m\) 的形式（某个单位 \(u \in R\)、某个 \(m \geq 0\)）。由此 \(R\) 的每个非零分式理想都具有形式 \(t^NR\)（某个 \(N \in \mathbb{Z}\)），从而是主分式理想，故由上一命题可逆。

反之，设 \(R\) 的每个非零分式理想都可逆。则由命题 9（3），\(R\) 的每个非零理想都有限生成，故 \(R\) 是诺特环。设 \(M\) 是 \(R\) 的唯一极大理想。若 \(M = M^2\)，则由 Nakayama 引理 \(M = 0\)，此时 \(R\) 是域，与假设矛盾。故存在 \(t \in M-M^2\). 由假设 \(M\) 可逆，而由于 \(t \in M\)，分式理想 \(tM^{-1}\) 是 \(R\) 中的非零理想。若 \(tM^{-1} \subseteq M\)，则 \(t \in M^2\)，与 \(t\) 的取法矛盾。故 \(tM^{-1} = R\)，即 \((t) = M\)，而 \(M\) 是非零主理想。由定理 7 的等价条件 4，\(R\) 是 D.V.R.，证明完毕。
\end{proof}

本节最后给出一个代数几何中的应用。

\noindent\textbf{仿射平面曲线的非奇异性与局部环}

设 \(k\) 是代数闭域，\(C\) 是 \(k\) 上的不可约仿射曲线。换句话说，\(C\) 是仿射代数集，其坐标环 \(k[C]\) 是整环，且其有理函数域 \(k(C)\) 在 \(k\) 上的超越次数为 1（参见第 15.4 节）。

回忆 \(C\) 上的点 \(v\) 按定义非奇异，若 \(\mathfrak{m}_{v,C}/\mathfrak{m}_{v,C}^2\) 是 \(k\) 上的 1 维向量空间，其中 \(\mathfrak{m}_{v,C}\) 是在 \(v\) 处有定义的有理函数的局部环 \(\mathcal{O}_{v,C}\) 中的唯一极大理想。

\begin{proposition}\label{pro:16.12}
设 \(v\) 是 \(k\) 上不可约仿射曲线 \(C\) 上的点。则 \(C\) 在 \(v\) 处非奇异当且仅当局部环 \(\mathcal{O}_{v,C}\) 是离散赋值环。
\end{proposition}

\begin{proof}
先设 \(v\) 非奇异。则 \(\dim_k(\mathfrak{m}_{v,C}/\mathfrak{m}_{v,C}^2) = 1\)，而由于 \(\mathcal{O}_{v,C}\) 是诺特环，由第 1 节习题 12 可知 \(\mathfrak{m}_{v,C}\) 是主理想。故由定理 7（4），\(\mathcal{O}_{v,C}\) 是 D.V.R. 反之，设 \(\mathcal{O}_{v,C}\) 是 D.V.R.，\(t\) 是 \(\mathcal{O}_{v,C}\) 的一致化元素。则 \(\mathfrak{m}_{v,C}\) 中每个元素都可以唯一地写成 \(at\) 的形式（某个 \(a \in \mathcal{O}_{v,C}\)）。把 \(at\) 映到 \(a \bmod \mathfrak{m}_{v,C}\) 的映射 \(\mathfrak{m}_{v,C} \to \mathcal{O}_{v,C}/\mathfrak{m}_{v,C}\) 容易验证是核为 \(\mathfrak{m}_{v,C}^2\) 的满 \(\mathcal{O}_{v,C}\)-模同态。故作为 \(\mathcal{O}_{v,C}/\mathfrak{m}_{v,C}\)-模有 \(\mathfrak{m}_{v,C}/\mathfrak{m}_{v,C}^2 \cong \mathcal{O}_{v,C}/\mathfrak{m}_{v,C}\). 由于 \(\mathcal{O}_{v,C}/\mathfrak{m}_{v,C} \cong k\)（第 15.4 节命题 46（5）），可知 \(\dim_k(\mathfrak{m}_{v,C}/\mathfrak{m}_{v,C}^2) = 1\)，故 \(v\) 是 \(C\) 上的非奇异点。
\end{proof}

\begin{definition}
若 \(v\) 是 \(C\) 上的非奇异点，相应的离散赋值 \(v_v\) 定义在 \(k(C)\) 上，则对 \(f \in k(C)\)，\(v_v(f) = n\) 是 \(f\) 在 \(v\) 处的\textbf{零点阶}（若 \(n \geq 0\)）或 \(f\) 在 \(v\) 处的\textbf{极点阶}（若 \(n < 0\)）。
\end{definition}

用命题 12 中曲线上点非奇异性的判别法，我们可以证明第 15.4 节中首次提到的结果：

\begin{corollary}\label{cor:16.13}
代数闭域 \(k\) 上的不可约仿射曲线 \(C\) 光滑当且仅当其坐标环 \(k[C]\) 整闭。
\end{corollary}

\begin{proof}
曲线 \(C\) 光滑当且仅当每个局部化 \(\mathcal{O}_{v,C}\) 都是 D.V.R. 由于 \(k[C]\) 的 Krull 维数为 1（第 1 节习题 11），每个 \(\mathcal{O}_{v,C}\) 也是如此。于是由定理 7（5），每个局部化 \(\mathcal{O}_{v,C}\) 都是 D.V.R. 当且仅当 \(\mathcal{O}_{v,C}\) 整闭。由第 15.4 节命题 49，这又等价于 \(k[C]\) 整闭，推论得证。
\end{proof}

\subsection*{习题}

\begin{enumerate}
\item 设 \(R\) 是关于其分式域 \(K\) 上赋值 \(v\) 的离散赋值环。若 \(x, y \in K\) 且 \(v(x) < v(y)\)，证明 \(v(x+y) = \min(v(x), v(y))\).〔注意 \(x+y = x(1+y/x)\).〕

\item 设 \(R\) 是唯一极大理想为 \(M\)、商环为 \(F = R/M\) 的离散赋值环。对任意 \(n \geq 0\) 证明 \(M^n/M^{n+1}\) 是 \(F\) 上的向量空间，且 \(\dim_F(M^n/M^{n+1}) = 1\).

\item 设 \(R\) 是整环，且是唯一极大理想 \(M = (t)\) 非零且为原理想的局部环，并假设 \(\bigcap_{n \geq 1}(t^n) = 0\). 证明 \(R\) 是离散赋值环。〔证明 \(R\) 中每个非零元素都可以写成 \(ut^n\) 的形式（某个单位 \(u \in R\)、某个 \(n \geq 0\)）。〕

\item 设 \(R\) 是唯一极大理想 \(M = (t)\) 为原理想的诺特局部环。证明或者 \(R\) 是离散赋值环，或者对某个 \(n \geq 0\) 有 \(t^n = 0\). 在后一情形证明 \(R\) 是 Artin 环。

\item 设 \(R\) 是诺特整环，且是 Krull 维数为 1 的局部环。设 \(M\) 是 \(R\) 的唯一极大理想，\(F = R/M\)，故 \(M/M^2\) 是 \(F\) 上的向量空间。

（a）证明若 \(\dim_F(M/M^2) = 1\) 则 \(R\) 是离散赋值环。

（b）若 \(R\) 的每个非零理想都是 \(M\) 的幂，证明 \(R\) 是离散赋值环。

\item 设 \(R\) 是分式域为 \(K\) 的整环。证明 \(K\) 的每个有限生成 \(R\)-子模都是 \(R\) 的分式理想。若 \(R\) 是诺特环，证明 \(A\) 是 \(R\) 的分式理想当且仅当 \(A\) 是 \(K\) 的有限生成 \(R\)-子模。

\item 若 \(R\) 是整环、\(A\) 是 \(R\) 的分式理想，证明若 \(A\) 投射则 \(A\) 有限生成。由此证明每个非诺特整环都含有不投射的理想。

\item 设 \(R\) 是诺特整环，且是极大理想 \(M\) 非零的局部环。证明 \(R\) 是 D.V.R. 当且仅当 \(R\) 中唯一的 \(M\)-准素理想是 \(M\) 的幂。

\item 设 \(C = Z(xz-y^2,\ yz-x^3,\ z^2-x^2y) \subset \mathbb{A}^3\)（代数闭域 \(k\) 上）。若 \(v = (0,0,0) \in C\)，证明 \(\dim_k(\mathfrak{m}_{v,C}/\mathfrak{m}_{v,C}^2) = 3\)，故 \(v\) 是 \(C\) 上的奇异点。由此证明 \(k[C]\) 在 \(k(C)\) 中不整闭，并确定它的整闭包。〔参见第 15.4 节习题 27。〕
\end{enumerate}

\section{Dedekind 整环}

上一节我们证明了离散赋值环就是 Krull 维数为 1 的整闭诺特整环中的局部环。本节我们考虑放松「局部环」这一条件的效果：

\begin{definition}
\textbf{Dedekind 整环}是 Krull 维数为 1 的诺特、整闭整环。等价地，\(R\) 是 Dedekind 整环，若 \(R\) 是诺特、整闭的整环、不是域，且其中每个非零素理想都是极大理想。
\end{definition}

第一个结果表明 Dedekind 整环是主理想整环这一类的推广。我们将在后面看到（定理 22）Dedekind 整环上有限生成模的结构定理，它推广了第 12.1 节中证明的主理想整环上的相应结果。

\begin{proposition}\label{pro:16.14}
（1）每个主理想整环都是 Dedekind 整环。

（2）代数数域的整数环是 Dedekind 整环。
\end{proposition}

\begin{proof}
主理想整环显然是诺特环，由于它是唯一分解整环而整闭（第 15.3 节例 3），且非零素理想都是极大理想（第 8.2 节命题 7），这证明了（1）。设 \(\mathcal{O}_K\) 是数域 \(K\) 的整数环，即 \(\mathbb{Z}\) 在 \(K\) 中的整闭包。则第 15.3 节推论 25 表明 \(\mathcal{O}_K\) 整闭，第 15.3 节定理 29 表明 \(\mathcal{O}_K\) 是诺特环，而 \(\mathcal{O}_K\) 中非零素理想都是极大理想这一事实已在同一定理之后的讨论中证明。这证明了（2）。
\end{proof}

下面的定理给出 Dedekind 整环的若干重要的等价刻画。回忆分式理想的基本性质已在上一节中发展。

\begin{theorem}\label{thm:16.15}
设 \(R\) 是分式域 \(K \neq R\) 的整环。下述条件是 \(R\) 为 Dedekind 整环的等价刻画：

（1）\(R\) 是诺特、整闭的环，且每个非零素理想都是极大理想；

（2）\(R\) 是诺特环，且对 \(R\) 的每个非零素理想 \(P\)，局部化 \(R_P\) 是离散赋值环；

（3）\(K\) 中 \(R\) 的每个非零分式理想都可逆；

（4）\(K\) 中 \(R\) 的每个非零分式理想都是投射 \(R\)-模；

（5）\(R\) 的每个非零真理想 \(I\) 都可以写成素理想的有限乘积：\(I = P_1P_2\cdots P_n\)（不必互不相同）。

当（5）中的条件成立时，素理想集合 \(\{P_1, \ldots, P_n\}\) 唯一确定，故 \(R\) 的每个非零真理想 \(I\) 都可以唯一地（不计次序）写成素理想幂之积。
\end{theorem}

\begin{proof}
若 \(R\) 满足（1），则由推论 8 \(R_P\) 是 D.V.R.，故（1）蕴涵（2）。反之，设每个 \(R_P\) 都是 D.V.R. 则由第 15.4 节命题 49，\(R\) 整闭；由第 15.4 节命题 46（3），每个非零素理想都是极大理想，故（2）蕴涵（1）。

现在设（1）成立，\(A\) 是 \(R\) 的非零分式理想。如命题 9 中那样令 \(A' = \{x \in K \mid xA \subseteq R\}\). 对 \(R\) 的任何素理想 \(P\)，\(R\)-模在局部化下的行为表明 \((AA')_P = A_P(A')_P = A_P(A_P)'\)（参见习题 4）。由于由已证部分 \(R_P\) 是 D.V.R.，由命题 9 有 \(A_P(A_P)' = R_P\). 于是对 \(R\) 的所有非零素理想 \(P\) 有 \((AA')_P = R_P\)，故由第 15.4 节习题 13 得 \(AA' = R\)，即 \(A\) 可逆，这表明（1）蕴涵（3）。反之，设 \(R\) 的每个非零分式理想都可逆。则由命题 9（3），\(R\) 中每个理想都有限生成，故 \(R\) 是诺特环。\(R\) 在非零素理想 \(P\) 处的每个局部化 \(R_P\) 都是其中非零分式理想可逆的局部环（参见习题 4），故由命题 11 是 D.V.R. 于是（3）蕴涵（2），从而（1）、（2）、（3）等价。它们与（4）的等价性由命题 10 给出。

现在设（1）成立，\(I\) 是 \(R\) 中任意非零真理想。由于 \(R\) 是诺特环，由第 15.2 节定理 21，\(I\) 有极小准素分解 \(I = Q_1 \cap \cdots \cap Q_n\). 相伴素理想 \(P_i = \operatorname{rad}Q_i\)（\(i = 1, \ldots, n\)）互不相同，而由于由假设 \(R\) 中素理想都是极大理想，这些相伴素理想两两互余，由此容易看出各 \(Q_i\) 也两两互余（习题 5）。于是由第 7.6 节定理 17，\(Q_1 \cap \cdots \cap Q_n = Q_1\cdots Q_n\)，故 \(I\) 是准素理想之积。\(R\) 的 \(P\)-准素理想与局部化 \(R_P\) 中的 \(PR_P\)-准素理想一一对应（第 15.4 节命题 42（3）），而由于 \(R_P\) 是 D.V.R.（因为（1）蕴涵（2）），由推论 6 可知若 \(Q\) 是 \(R\) 中的 \(P\)-准素理想，则对某个整数 \(m \geq 1\) 有 \(Q = P^m\). 把这一点应用于 \(Q_i\)（\(i = 1, \ldots, n\)）可知 \(I\) 是素理想幂之积，这给出（5）的第一个蕴涵。

反之，设 \(R\) 的所有非零真理想都可以写成素理想之积。我们先证明对任何整环，理想分解为可逆素理想之积是唯一的，即若 \(P_1\cdots P_n = P'_1\cdots P'_m\) 是 \(I\) 分解为可逆素理想之积的两种分解，则 \(n = m\) 且两个素理想集合 \(\{P_1, \ldots, P_n\}\) 与 \(\{P'_1, \ldots, P'_m\}\) 相同。设 \(P_1\) 是集合 \(\{P_1, \ldots, P_m\}\) 中的极小元。由于 \(P_1\cdots P_n \subseteq P_1\)，素理想 \(P_1\) 包含 \(P_1, \ldots, P_n\) 中的一个，比如 \(P_1 \subseteq P_1\). 类似地 \(P_1\) 包含某个 \(P'_j\)（\(j = 1, \ldots, m\)）。则 \(P'_j \subseteq P_1 \subseteq P_1\)，而由 \(P_1\) 的极小性得 \(P'_j = P_1 = P_1\)，故分解化为 \(P_1P_2\cdots P_n = P_1P'_2\cdots P'_m\). 由于 \(P_1\) 可逆，乘以逆理想表明 \(P_2\cdots P_n = P'_2\cdots P'_m\)，简单的归纳即可完成证明。特别地，由于 Dedekind 整环中每个分式理想（特别是 \(R\) 的每个素理想）都可逆，（5）中的唯一性论断现在由（5）的第一个论断推出。

下面我们证明 \(R\) 中可逆素理想都是极大理想。设 \(P\) 是 \(R\) 中的可逆素理想，取 \(a \in R\)、\(a \notin P\). 我们要证明 \(P+aR = R\). 由假设，两个理想 \(P+aR\) 与 \(P+a^2R\) 都可以写成素理想之积，比如 \(P+aR = P_1\cdots P_n\)、\(P+a^2R = P'_1\cdots P'_m\). 注意对 \(i = 1, \ldots, n\) 有 \(P \subseteq P_i\)，对 \(j = 1, \ldots, m\) 有 \(P \subseteq P'_j\). 在整环 \(R/P\) 中我们有分解 \((a) = (P_1/P)\cdots(P_n/P)\)，且每个 \(P_i/P\) 是 \(R/P\) 中的素理想。由于这个乘积是主理想，每个 \(P_i/P\) 也是可逆 \(R/P\)-理想（参见习题 2）。类似地 \((a^2) = (P'_1/P)\cdots(P'_m/P)\) 是分解为可逆素理想之积。于是 \((a)^2 = (P_1/P)^2\cdots(P_n/P)^2 = (P'_1/P)\cdots(P'_m/P)\) 给出整环 \(R/P\) 中分解为可逆素理想之积的两种分解，故由上一段的唯一性结果，\(m = 2n\) 且
\[
\{P_1/P, P_1/P, \ldots, P_n/P, P_n/P\} = \{P'_1/P, \ldots, P'_m/P\}.
\]
由此 \(R\) 中素理想集合 \(\{P'_1, \ldots, P'_m\}\) 由 \(P_1, \ldots, P_n\) 各重复两次组成。这表明 \(P+a^2R = (P+aR)^2\). 由于 \(P \subseteq P+a^2R\) 且 \((P+aR)^2 \subseteq P^2+aR\)，我们有 \(P \subseteq P^2+aR\)，故 \(P\) 中每个元素 \(x\) 都可以写成 \(x = y+az\)（\(y \in P^2\)、\(z \in R\)）。则 \(az = x-y \in P\)，而由于 \(a \notin P\)，得 \(z \in P\)，这表明 \(P \subseteq P^2+aP\). 显然 \(P^2+aP \subseteq P\)，故 \(P = P^2+aP = P(P+aR)\). 由于假设 \(P\) 可逆，可知对任何 \(a \in R-P\) 有 \(R = P+aR\)，这证明 \(P\) 是极大理想。

现在我们证明每个非零素理想都可逆。若 \(P\) 是非零素理想，取 \(P\) 中任意非零元素 \(a\). 由假设，\(Ra = P_1\cdots P_n\) 可以写成素理想之积，且由于它们的乘积是主理想，各 \(P_1, \ldots, P_n\) 可逆（再次利用习题 2）。由于 \(P_1\cdots P_n = Ra \subseteq P\)，素理想 \(P\) 包含某个 \(P_i\)（\(1 \leq i \leq n\)）。由上一段 \(P_i\) 是极大理想，故 \(P = P_i\) 可逆。

最后，由于 \(R\) 的每个非零真理想都是素理想之积，可知 \(R\) 的每个非零理想都可逆；又由于 \(R\) 的每个分式理想都具有形式 \((d^{-1})I\)（\(I\) 是 \(R\) 中的理想），\(R\) 的每个分式理想也都可逆。这证明了（5）蕴涵（3），从而完成了定理的证明。
\end{proof}

下述推论由命题 14 立即推出：

\begin{corollary}\label{cor:16.16}
若 \(\mathcal{O}_K\) 是代数数域 \(K\) 的整数环，则 \(\mathcal{O}_K\) 中每个非零理想 \(I\) 都可以唯一地写成互不相同的素理想幂之积。
\end{corollary}

\noindent\textbf{注：} 这里给出的 Dedekind 整环的展开顺序与历史发展相反。如第 9.3 节所述，非零理想分解为素理想之积的唯一性取代了数域整数环中非零元素分解为素元素之积的唯一性的失效。推论 16 中整数环的这一性质正是最初引出理想这一概念的原因，而 Dedekind 最初用定理 15 中的性质 5 定义我们现在称为 Dedekind 整环的对象。是 Noether 注意到它们也可以用性质（1）刻画，而我们就是把（1）作为 Dedekind 整环的最初定义。

Dedekind 整环中分解为素理想之积的唯一性可以用来显式定义 \(R\) 上的赋值，使得赋值环是定理 15（2）中的局部化 \(R_P\)（参见习题 6）。下面我们说明理想的唯一分解如何用来定义与 \(\mathbb{Z}\) 中整数整除性类似的理想的整除理论。

\begin{definition}
若 \(A\) 与 \(B\) 是整环 \(R\) 中的理想，则称 \(B\) \textbf{整除} \(A\)（且 \(A\) 被 \(B\) 整除），若存在 \(R\) 中的理想 \(C\) 使 \(A = BC\).
\end{definition}

若 \(B\) 整除 \(A\) 则当然 \(A \subseteq B\). 若 \(R\) 是 Dedekind 整环，则其逆命题也成立：\(A \subseteq B\) 蕴涵 \(C = AB^{-1} \subseteq BB^{-1} = R\)，故 \(C\) 是 \(R\) 中满足 \(BC = A\) 的理想。

我们也可以定义两个理想 \(A\) 与 \(B\) 的\textbf{最大公因子} \((A,B)\)：\((A,B)\) 整除 \(A\) 与 \(B\)，且任何同时整除 \(A\) 与 \(B\) 的理想都整除 \((A,B)\). 下一命题的第二个论断表明这个最大公因子对 Dedekind 整环中的整理想总是存在，并给出与两个整数的最大公因子公式类似的公式。

\begin{proposition}\label{pro:16.17}
设 \(R\) 是 Dedekind 整环，\(A, B\) 是 \(R\) 中两个非零理想，其素理想分解为 \(A = P_1^{e_1}\cdots P_n^{e_n}\)、\(B = P_1^{f_1}\cdots P_n^{f_n}\)（其中 \(e_i, f_i \geq 0\)，\(i = 1, \ldots, n\)）。则

（1）\(A \subseteq B\) 当且仅当 \(B\) 整除 \(A\)（即「包含即整除」），当且仅当对 \(i = 1, \ldots, n\) 有 \(f_i \leq e_i\)；

（2）\(A+B = (A,B) = P_1^{\min(e_1,f_1)}\cdots P_n^{\min(e_n,f_n)}\)，故特别地 \(A\) 与 \(B\) 互素（即 \(A+B = R\)）当且仅当它们没有公共素理想因子。
\end{proposition}

\begin{proof}
（1）的第一个论断已在上面证明。若每个 \(f_i \leq e_i\)，则取 \(C = P_1^{e_1-f_1}\cdots P_n^{e_n-f_n} \subseteq R\) 表明 \(B\) 整除 \(A\). 反之，若 \(B\) 整除 \(A\)，把 \(C\) 写成素理想之积并代入 \(A = BC\) 可知对所有 \(i\) 有 \(f_i \leq e_i\)，这证明了（1）的全部内容。由于 \(A+B\) 是同时包含 \(A\) 与 \(B\) 的最小理想，（2）现在由（1）推出。
\end{proof}

\begin{proposition}\label{pro:16.18}
（\textbf{中国剩余定理}）设 \(R\) 是 Dedekind 整环，\(P_1, P_2, \ldots, P_n\) 是 \(R\) 中互不相同的素理想，\(a_i \geq 0\) 是整数（\(i = 1, \ldots, n\)）。则
\[
R/(P_1^{a_1}P_2^{a_2}\cdots P_n^{a_n}) \cong R/P_1^{a_1} \times R/P_2^{a_2} \times \cdots \times R/P_n^{a_n}.
\]
等价地，对任意元素 \(r_1, r_2, \ldots, r_n \in R\) 存在元素 \(r \in R\)（在相差 \(P_1^{a_1}\cdots P_n^{a_n}\) 中元素的意义下唯一）满足相应的同余条件。
\end{proposition}

\begin{proof}
这由第 7.6 节定理 17 立即推出，因为上一命题表明各 \(P_i^{a_i}\) 是两两互余的理想。
\end{proof}

\begin{corollary}\label{cor:16.19}
设 \(I\) 是 Dedekind 整环 \(R\) 中的理想。则

（1）存在与 \(I\) 互素的理想 \(J\) 使乘积 \(IJ = (a)\) 是主理想；

（2）若 \(I\) 非零，则商环 \(R/I\) 中每个理想都是主理想；等价地，若 \(I_1\) 是 \(R\) 中含 \(I\) 的理想，则对某个 \(b \in R\) 有 \(I_1 = I+Rb\)；

（3）\(R\) 中每个理想都可以由两个元素生成；事实上若 \(I\) 非零且 \(0 \neq a \in I\)，则对某个 \(b \in I\) 有 \(I = Ra+Rb\).
\end{corollary}

\begin{proof}
设 \(I = P_1^{e_1}\cdots P_n^{e_n}\) 是 \(I\) 在 \(R\) 中的素理想分解。对每个 \(i = 1, \ldots, n\)，取 \(r_i \in P_i^{e_i}-P_i^{e_i+1}\). 由上一命题，存在 \(a \in R\) 使对所有 \(i\) 有 \(a \equiv r_i \pmod{P_i^{e_i+1}}\). 于是对所有 \(i\) 有 \(a \in P_i^{e_i}-P_i^{e_i+1}\)，故由命题 17（1），(a) 的素理想分解中 \(P_i\) 的幂恰为 \(e_i\)：
\[
(a) = P_1^{e_1}\cdots P_n^{e_n}P_{n+1}^{e_{n+1}}\cdots P_m^{e_m},
\]
其中 \(P_{n+1}, \ldots, P_m\) 是与 \(P_1, \ldots, P_n\) 不同的素理想。令 \(J = P_{n+1}^{e_{n+1}}\cdots P_m^{e_m}\) 即得（1）。对（2），由中国剩余定理，只需证明当 \(P\) 是素理想时 \(R/P^m\) 中每个理想都是主理想，而这立即推出，因为 \(R/P^m \cong R_P/P^mR_P\) 且局部化 \(R_P\) 是主理想整环。最后，（3）由（2）取 \(I = Ra\) 推出。
\end{proof}

推论 19 的第一个论断表明存在与 \(I\) 互素、且落在 \(R\) 的类群中 \(I\) 的逆类里的整理想 \(J\). 甚至可以给 \(J\) 加一些额外条件，参见习题 11。

\begin{corollary}\label{cor:16.20}
若 \(R\) 是 Dedekind 整环，则 \(R\) 是主理想整环（即 \(R\) 的类数为 1）当且仅当 \(R\) 是唯一分解整环。
\end{corollary}

\begin{proof}
每个主理想整环都是唯一分解整环，故设 \(R\) 是唯一分解整环，\(P\) 是 \(R\) 中任意非零素理想。则由推论 19，对某些 \(a, b \in R\) 有 \(P = Ra+Rb\). 由于 \(a\) 的不可约因子之积是 \(P\) 中的元素，对 \(a\) 的某个不可约因子 \(a'\) 有 \((a') \subseteq P\)，于是由命题 17（1）\(P\) 整除 \((a')\) 在 \(R\) 中的分解。由此 \(P = (a')\) 是主理想，因为 \((a')\) 是素理想（第 8.3 节命题 12）。由于 \(R\) 中每个理想都是素理想之积，\(R\) 的每个理想都是主理想，即 \(R\) 是主理想整环。
\end{proof}

推论 20 表明 Dedekind 整环 \(R\) 的类数给出了元素唯一分解失效程度的一个度量。代数数论的一个基本结果是代数数域整数环的类数有限。然而对一般的 Dedekind 整环，类数不必有限。事实上对任何阿贝尔群 \(A\)（有限或无限），都存在类群同构于 \(A\) 的 Dedekind 整环。

\noindent\textbf{Dedekind 整环上的模与有限生成模基本定理}

下面我们转向 Dedekind 整环 \(R\) 上的模。\(R\) 的每个分式理想都是 \(R\)-模，而下一命题的第一个论断表明 \(R\) 的两个分式理想作为 \(R\)-模同构当且仅当它们代表 \(R\) 的类群中的同一个元素。

\begin{proposition}\label{pro:16.21}
设 \(R\) 是分式域为 \(K\) 的 Dedekind 整环。

（1）设 \(I\) 与 \(I'\) 是 \(R\) 的两个分式理想。则作为 \(R\)-模有 \(I \cong I'\) 当且仅当 \(I\) 与 \(I'\) 相差一个非零主分式理想：对某个 \(0 \neq a \in K\) 有 \(I' = (a)I\).

（2）更一般地，设 \(I_1, I_2, \ldots, I_n\) 与 \(I'_1, I'_2, \ldots, I'_m\) 是 Dedekind 整环 \(R\) 的分式域 \(K\) 中的非零分式理想。则作为 \(R\)-模有
\[
I_1 \oplus I_2 \oplus \cdots \oplus I_n \cong I'_1 \oplus I'_2 \oplus \cdots \oplus I'_m
\]
当且仅当 \(n = m\) 且乘积理想 \(I_1I_2\cdots I_n\) 与 \(I'_1I'_2\cdots I'_m\) 相差一个主理想：
\[
I_1I_2\cdots I_n = (a)I'_1I'_2\cdots I'_m
\]
（某个 \(0 \neq a \in K\)）。

（3）特别地，
\[
I_1 \oplus I_2 \oplus \cdots \oplus I_n \cong R \oplus \cdots \oplus R \oplus (I_1I_2\cdots I_n) \quad (\text{\(n-1\) 个 \(R\)}),
\]
且 \(R \oplus I \cong R \oplus I'\) 当且仅当 \(I\) 与 \(I'\) 相差一个主理想：\(I' = (a)I\)（某个 \(a \in K\)）。
\end{proposition}

\begin{proof}
乘以 \(0 \neq a \in K\) 给出从 \(I\) 到 \((a)I\) 的 \(R\)-模同构，故若 \(I' = (a)I\) 则作为 \(R\)-模有 \(I \cong I'\). 反之，注意我们可以假设 \(I' \neq 0\)，此时 \(I \cong I'\) 蕴涵 \(R \cong I'^{-1}I\). 但这说明 \(I'^{-1}I = aR\) 是主理想（其生成元 \(a\) 由 \(1 \in R\) 的像给出），即 \(I = (a)I'\)，这证明了（1）。

下面我们证明对任何非零分式理想 \(I\) 与 \(J\) 有 \(I \oplus J \cong R \oplus IJ\). 必要时用同构的 \(R\)-模 \(aI\) 与 \(bJ\) 代替 \(I\) 与 \(J\)，我们可以假设 \(I\) 与 \(J\) 是互素的整理想（参见习题 12），故 \(I+J = R\) 且 \(I \cap J = IJ\). 容易看出把 \((x,y)\) 映到 \(x+y\) 的映射 \(I \oplus J \to I+J = R\) 是满 \(R\)-模同态，其核为 \(IJ\)，故有 \(R\)-模的正合序列
\[
0 \to IJ \to I \oplus J \to R \to 0.
\]
由于 \(R\) 是自由模，这个序列分裂，故 \(I \oplus J \cong R \oplus IJ\)，正如所断言。（3）的第一个论断现在由归纳推出，而把这个论断与（1）结合可知若对某个非零 \(a \in K\) 有 \(I_1\cdots I_n = (a)I'_1\cdots I'_n\)，则 \(I_1 \oplus \cdots \oplus I_n\) 同构于 \(I'_1 \oplus \cdots \oplus I'_n\). 这证明了（2）中「若」的部分。还需证明（2）中「仅当」的部分，因为（3）中的相应论断是其特殊情形。

于是设作为 \(R\)-模有 \(I_1 \oplus I_2 \oplus \cdots \oplus I_n \cong I'_1 \oplus I'_2 \oplus \cdots \oplus I'_m\). 由于 \(I \otimes_RK\) 是理想 \(I\) 在 \(K\) 中的局部化（参见第 15.4 节命题 41），可知对 \(K\) 的任何非零分式理想 \(I\) 有 \(I \otimes_RK \cong K\). 由于张量积与直和可交换，\((I_1 \oplus \cdots \oplus I_n)\otimes_RK \cong K^n\) 是 \(K\) 上的 \(n\) 维向量空间。类似地 \(I'_1 \oplus \cdots \oplus I'_m \otimes_RK \cong K^m\)，由此 \(n = m\).

注意对任何非零元素 \(a_1 \in I_1\)，用同构的分式理想 \(a_1^{-1}I_1\) 代替 \(I_1\) 不影响（2）中各论断的正确性。故我们可以假设 \(I_1\) 包含 \(R\)，类似地可以假设（2）中每个分式理想都包含 \(R\). 用 \(\varphi\) 表示从 \(I_1 \oplus \cdots \oplus I_n\) 到 \(I'_1 \oplus \cdots \oplus I'_n\) 的 \(R\)-模同构。对 \(i = 1, 2, \ldots, n\) 定义
\[
\varphi((0, \ldots, 0, 1, 0, \ldots, 0)) = (a_{1,i}, a_{2,i}, \ldots, a_{n,i}) \in I'_1 \oplus I'_2 \oplus \cdots \oplus I'_n,
\]
其中左边的 \(1 \in I_i\) 位于第 \(i\) 个位置。由于 \(\varphi\) 是 \(R\)-模同态，可知对每个 \(j = 1, 2, \ldots, n\) 有
\[
I'_j = a_{j,1}I_1+a_{j,2}I_2+\cdots+a_{j,j}I_j+\cdots+a_{j,n}I_n.
\]
对这些理想取乘积（\(j = 1, 2, \ldots, n\)）可知对 \(\{1, 2, \ldots, n\}\) 的任何置换 \(\{j_1, j_2, \ldots, j_n\}\) 有
\[
a_{1,j_1}a_{2,j_2}\cdots a_{n,j_n}I_1I_2\cdots I_n \subseteq I'_1I'_2\cdots I'_n.
\]
于是
\[
d\,I_1I_2\cdots I_n \subseteq I'_1I'_2\cdots I'_n,
\]
其中 \(d\) 是矩阵 \((a_{i,j})\) 的行列式，因为行列式是形如 \(\varepsilon(\sigma)a_{1,\sigma(1)}\cdots a_{n,\sigma(n)}\) 的项之和（\(\varepsilon(\sigma)\) 是 \(\{1,2,\ldots,n\}\) 的置换 \(\sigma\) 的符号）。类似地，对 \(j = 1, \ldots, n\) 定义
\[
\varphi^{-1}((0, \ldots, 0, 1, 0, \ldots, 0)) = (b_{1,j}, b_{2,j}, \ldots, b_{n,j}) \in I_1 \oplus I_2 \oplus \cdots \oplus I_n,
\]
其中左边的 \(1 \in I'_j\) 位于第 \(j\) 个位置。两个矩阵 \((a_{i,j})\) 与 \((b_{i,j})\) 的乘积就是单位矩阵，故 \(d \neq 0\) 且矩阵 \((b_{i,j})\) 的行列式是 \(d^{-1}\). 如上有
\[
d^{-1}I'_1I'_2\cdots I'_n \subseteq I_1I_2\cdots I_n,
\]
这表明对某个 \(0 \neq a = d \in K\) 有 \(I_1I_2\cdots I_n = (a)I'_1I'_2\cdots I'_n\)，命题证明完毕。
\end{proof}

下面我们考虑 Dedekind 整环上有限生成模，并证明推广第 12 章中主理想整环上有限生成模结果的结构定理。

回忆 \(M\) 的\textbf{秩}是 \(M\) 中 \(R\)-线性无关元素的最大个数，等价地，是 \(M \otimes_RK\) 作为 \(K\)-向量空间的维数，其中 \(K\) 是 \(R\) 的分式域（参见第 12.1 节习题 1--4、20）。

\begin{theorem}\label{thm:16.22}
设 \(M\) 是 Dedekind 整环 \(R\) 上的有限生成模。用 \(n \geq 0\) 表示 \(M\) 的秩，\(\operatorname{Tor}(M)\) 表示 \(M\) 的挠子模。则对 \(R\) 的某个理想 \(I\) 有
\[
M \cong \underbrace{R \oplus R \oplus \cdots \oplus R}_{n \text{ 个}} \oplus I \oplus \operatorname{Tor}(M),
\]
且
\[
\operatorname{Tor}(M) \cong R/P_1^{e_1} \times R/P_2^{e_2} \times \cdots \times R/P_s^{e_s}
\]
（某个 \(s \geq 0\)，各 \(P_i^{e_i}\) 是（不必互不相同的）素理想的幂，\(e_i \geq 1\)）。理想 \(P_i\)（\(i = 1, \ldots, s\)）唯一，理想 \(I\) 在相差一个主理想的乘法的意义下唯一。
\end{theorem}

\begin{proof}
先设 \(M\) 是 \(R\) 上的有限生成无挠模，即 \(\operatorname{Tor}(M) = 0\). 则从 \(M\) 到 \(M \otimes_RK\) 的自然 \(R\)-模同态是单射，故可以把 \(M\) 看作向量空间 \(M \otimes_RK\) 的 \(R\)-子模。若 \(M\) 在 \(R\) 上的秩为 \(n\)，则 \(M \otimes_RK\) 是 \(K\) 上 \(n\) 维向量空间。设 \(x_1, \ldots, x_n\) 是 \(M \otimes_RK\) 在 \(K\) 上的基，\(m_1, \ldots, m_s\) 是 \(M\) 的 \(R\)-模生成元。每个 \(m_i\)（\(i = 1, \ldots, s\)）都可以写成 \(x_1, \ldots, x_n\) 的 \(K\)-线性组合。设 \(0 \neq d \in R\) 是这些线性组合中所有 \(K\)-系数的一个公分母，令 \(y_i = x_i/d\)（\(i = 1, \ldots, n\)）。则
\[
M \subseteq Ry_1+\cdots+Ry_n \subseteq Kx_1+\cdots+Kx_n,
\]
这表明 \(M\) 包含在秩为 \(n\) 的自由 \(R\)-子模中，且 \(M\) 中每个元素都可以唯一地写成
\[
m = a_1y_1+\cdots+a_ny_n \qquad (a_1, \ldots, a_n \in R)
\]
的形式。由 \(\varphi(a_1y_1+\cdots+a_ny_n) = a_n\) 定义的映射 \(\varphi : M \to R\) 是 \(R\)-模同态，故有正合序列
\[
0 \to \ker\varphi \to M \xrightarrow{\varphi} I_1 \to 0,
\]
其中 \(I_1\) 是 \(\varphi\) 在 \(R\) 中的像，从而是 \(R\) 中的理想。子模 \(\ker\varphi\) 也是无挠 \(R\)-模，其秩至多为 \(n-1\)（因为它包含在 \(Ry_1+\cdots+Ry_{n-1}\) 中），比较秩可知 \(I_1\) 非零且 \(\ker\varphi\) 的秩恰为 \(n-1\). 由定理 15（4），\(I_1\) 是投射 \(R\)-模，故这个序列分裂：
\[
M \cong I_1 \oplus (\ker\varphi).
\]
对秩归纳可知有限生成无挠 \(R\)-模同构于 \(n\) 个非零理想的直和：
\[
M \cong I_1 \oplus I_2 \oplus \cdots \oplus I_n.
\]
由于 \(I_1, \ldots, I_n\) 都是投射 \(R\)-模，可知任何有限生成无挠 \(R\)-模都是投射模。

现在设 \(M\) 是任意有限生成 \(R\)-模。商 \(M/\operatorname{Tor}(M)\) 有限生成且无挠，从而由刚才证明的部分是投射模。因此正合序列
\[
0 \to \operatorname{Tor}(M) \to M \to M/\operatorname{Tor}(M) \to 0
\]
分裂，故
\[
M \cong \operatorname{Tor}(M) \oplus (M/\operatorname{Tor}(M)).
\]
由上一段的结果，\(M/\operatorname{Tor}(M)\) 同构于 \(R\) 的 \(n\) 个非零理想的直和，而由命题 21 我们得到对 \(R\) 的某个理想 \(I\) 有
\[
M \cong \underbrace{R \oplus R \oplus \cdots \oplus R}_{n \text{ 个}} \oplus I \oplus \operatorname{Tor}(M).
\]
关于理想 \(I\) 的唯一性论断也由命题 21（3）中的唯一性论断立即推出。

还需证明关于挠子模 \(\operatorname{Tor}(M)\) 的论断。设 \(N\) 是有限生成挠 \(R\)-模。设 \(I = \operatorname{Ann}(N)\) 是 \(N\) 在 \(R\) 中的零化子，并设 \(I = P_1^{e_1}\cdots P_s^{e_s}\) 是 \(I\) 在 \(R\) 中的素理想分解，其中 \(P_1, \ldots, P_s\) 是互不相同的素理想。则 \(N\) 是 \(R/I\)-模，且
\[
R/I \cong R/P_1^{e_1} \times R/P_2^{e_2} \times \cdots \times R/P_s^{e_s}.
\]
由此作为 \(R\)-模有
\[
N \cong (N/P_1^{e_1}N) \times (N/P_2^{e_2}N) \times \cdots \times (N/P_s^{e_s}N).
\]
每个 \(N/P_i^{e_i}N\) 是 \(R/P_i^{e_i} \cong R_{P_i}/P_i^{e_i}R_{P_i}\) 上的有限生成模，其中 \(R_{P_i}\) 是 \(R\) 在素理想 \(P_i\) 处的局部化，即是被 \(P_i^{e_i}R_{P_i}\) 零化的 \(R_{P_i}\) 上的有限生成模。由于 \(R\) 是 Dedekind 整环，每个 \(R_{P_i}\) 是主理想整环（甚至是 D.V.R.），故我们可以应用主理想整环上有限生成模的基本定理，得出每个 \(N/P_i^{e_i}N\) 作为 \(R_{P_i}\)-模同构于有限多个形如 \(R_{P_i}/P_i^fR_{P_i}\)（\(f \leq e_i\)）的模的直和。由此每个 \(N/P_i^{e_i}N\) 作为 \(R\)-模同构于有限多个形如 \(R/P_i^fR\)（\(f \leq e_i\)）的模的直和。这证明 \(N\) 同构于有限多个形如 \(R/P_i^f\)（各 \(P_i\) 为不同素理想）的模的直和。故 \(\operatorname{Tor}(M)\) 可以如定理中那样分解为直和。

最后，还需证明 \(\operatorname{Tor}(M)\) 的分解中各理想 \(P_i^{e_i}\)（\(i = 1, \ldots, s\)）唯一。这与第 12.1 节定理 10 证明中的唯一性论证类似（另见第 12.1 节习题 11--12）：对 \(R\) 的任何素理想 \(P\)，商 \(P^{i-1}M/P^iM\) 是域 \(R/P\) 上的向量空间，而差 \(\dim_{R/P}P^{i-1}M/P^iM-\dim_{R/P}P^iM/P^{i+1}M\) 是同构于 \(R/P^i\) 的直和项的个数，从而由 \(M\) 唯一确定。定理证明完毕。
\end{proof}

若 \(M\) 是 Dedekind 整环 \(R\) 上的有限生成模（如定理 22 中那样），则 \(M\) 作为 \(R\)-模的同构型由秩 \(n\)、素幂 \(P_i^{e_i}\)（\(i = 1, \ldots, s\)，称为 \(M\) 的\textbf{初等因子}）以及理想 \(I\) 在 \(R\) 的类群中的类（称为 \(M\) 的 \textbf{Steinitz 类}）决定。注意主理想整环就是类数为 1 的 Dedekind 整环，此时 \(R\) 的每个非零理想 \(I\) 作为 \(R\)-模都同构于 \(R\). 在这种情形，定理 22 归结为第 12 章中主理想整环上有限生成模结构定理的初等因子形式。定理 22 中对挠 \(R\)-模的描述也有不变因子形式（参见习题 14）。

下一个结果把主理想整环上有限生成投射模的刻画（第 12.1 节习题 21）推广到 Dedekind 整环。

\begin{corollary}\label{cor:16.23}
Dedekind 整环上的有限生成模是投射模当且仅当它是无挠模。
\end{corollary}

\begin{proof}
我们在定理 22 的证明中已经证明有限生成无挠 \(R\)-模是投射模，故由定理 22 中 \(M\) 的分解，\(M\) 投射当且仅当 \(\operatorname{Tor}(M)\) 投射（参见第 10.5 节习题 3）。为完成证明，只需证明非零挠 \(R\)-模不投射，这留作习题（参见习题 15）。
\end{proof}

\subsection*{习题}

\begin{enumerate}
\item 若 \(R\) 是整环，证明 \(R\) 的每个分式理想都可逆当且仅当 \(R\) 的每个整理想都可逆。

\item 设 \(R\) 是分式域为 \(K\) 的整环，\(A_1, A_2, \ldots, A_n\) 是 \(R\) 的分式理想，其乘积是非零主分式理想：\(A_1A_2\cdots A_n = Rx\)（某个 \(0 \neq x \in K\)）。对每个 \(i = 1, \ldots, n\) 证明 \(A_i\) 是可逆分式理想，其逆为 \((x^{-1})A_1\cdots A_{i-1}A_{i+1}\cdots A_n\).

\item 设 \(R\) 是分式域为 \(K\) 的整环，\(P\) 是 \(R\) 中的非零素理想。证明 \(R_P\) 在 \(K\) 中的分式理想是形如 \(AR_P\) 的 \(R_P\)-模，其中 \(A\) 是 \(R\) 的分式理想。

\item 设 \(R\) 是分式域为 \(K\) 的整环，\(A\) 是 \(K\) 中 \(R\) 的分式理想。如命题 9 中那样令 \(A' = \{x \in K \mid xA \subseteq R\}\).

（a）对 \(R\) 中任何素理想 \(P\)，证明 \(A'\) 在 \(P\) 处的局部化 \((A')_P\) 是 \(K\) 中 \(R_P\) 的分式理想。

（b）若 \(A\) 是有限生成 \(R\)-模，证明 \((A')_P = (A_P)'\)，其中 \((A_P)'\) 是对应于局部化 \(A_P\) 的 \(R_P\)-分式理想 \(\{x \in K \mid xA_P \subseteq R_P\}\).

\item 若 \(Q_1\) 是 \(P_1\)-准素理想、\(Q_2\) 是 \(P_2\)-准素理想，其中 \(P_1\) 与 \(P_2\) 是诺特环 \(R\) 中互余的理想，证明 \(Q_1\) 与 \(Q_2\) 也互余。〔利用第 15.2 节命题 14。〕

\item 设 \(R\) 是分式域为 \(K\) 的 Dedekind 整环。

（a）证明 \(K\) 中 \(R\) 的每个非零分式理想都可以唯一地写成互不相同的素幂之积 \(P_1^{a_1}\cdots P_n^{a_n}\)，其中各 \(a_i\) 是非零整数（可以为负）。

（b）若 \(0 \neq x \in K\)，用 \(v_P(x)\) 表示主理想 \((x)\) 的分解中素理想 \(P\) 的幂（如（a）中那样；若 \(P\) 不是出现的素理想之一则 \(v_P(x) = 0\)）。证明 \(v_P\) 是 \(K\) 上的赋值，其赋值环是 \(R_P\)，即 \(R\) 在 \(P\) 处的局部化。

\item 设 \(R\) 是诺特整环且不是域。证明 \(R\) 是 Dedekind 整环当且仅当对 \(R\) 的每个极大理想 \(M\)，都不存在满足 \(M \subsetneq I \subsetneq M^2\) 的 \(R\) 的理想 \(I\).〔利用第 1 节习题 12 以及定理 7 与 15。〕

\item 设 \(R\) 是 Krull 维数为 1 的诺特整环。证明 \(R\) 中每个非零理想 \(I\) 都可以唯一地写成根互不相同的准素理想之积。〔参见定理 15 的证明。利用极小准素分解中属于孤立素理想的准素分量的唯一性（第 15.2 节定理 21）。〕

\item 设 \(R\) 是整环。证明对每个非零素理想 \(P\)，\(R_P\) 是 D.V.R. 当且仅当对每个非零极大理想，\(R_M\) 是 D.V.R.

\item 设 \(R\) 是诺特整环且不是域。证明 \(R\) 是 Dedekind 整环当且仅当非零素理想 \(M\) 都是极大理想，且每个 \(M\)-准素理想都是 \(M\) 的幂。

\item 若 \(I\) 与 \(J\) 是 Dedekind 整环 \(R\) 中的非零理想，证明存在 \(R\) 中的整理想 \(I_1\)，它与 \(I\) 与 \(J\) 都互素且使 \(I_1I\) 是 \(R\) 中的主理想。

\item 若 \(I\) 与 \(J\) 是 Dedekind 整环 \(R\) 的非零分式理想，证明存在 \(a, \beta \in K\) 使 \(aI\) 与 \(\beta J\) 是 \(R\) 中互素的非零整理想。

\item 设 \(I\) 与 \(J\) 是 Dedekind 整环 \(R\) 中的非零理想。证明存在满足 \(I_1 \subseteq I\) 且与 \(J\) 互素的理想 \(I_1\).〔用推论 19 求理想 \(I_2\) 使 \(I_2I = (a)\) 且 \((I_2, J) = R\). 若 \(I_2 = P_1^{e_1}\cdots P_n^{e_n}\)，取 \(b \in R\) 使 \(b \in P_i^{e_i}-P_i^{e_i+1}\) 且对整除 \(J\) 的每个素理想 \(P\) 有 \(b \equiv 1 \pmod P\). 证明对某个理想 \(I_1\) 有 \((b) = I_2I_1\)，并考虑 \((a)I_1\) 以证明 \(I_1 \subseteq I\).〕

\item 证明 Dedekind 整环 \(R\) 上每个有限生成挠模都同构于直和 \(R/I_1 \oplus R/I_2 \oplus \cdots \oplus R/I_n\)，其中 \(I_1, \ldots, I_n\) 是 \(R\) 中满足 \(I_1 \subseteq I_2 \subseteq \cdots \subseteq I_n\) 的唯一非零理想（称为 \(M\) 的\textbf{不变因子}）。〔参见第 12.1 节。〕

\item 若 \(P\) 是 Dedekind 整环 \(R\) 中的非零素理想，证明对任何 \(n \geq 1\)，\(R/P^n\) 不是投射 \(R\)-模。〔考虑正合序列 \(0 \to P^n/P^{n+1} \to R/P^{n+1} \to R/P^n \to 0\).〕由此证明若 \(M \neq 0\) 是有限生成挠 \(R\)-模，则 \(M\) 不是投射模。〔参见第 10.5 节习题 3。〕

\item 证明 Dedekind 整环 \(R\) 的类数为 1 当且仅当每个有限生成投射 \(R\)-模都是自由模。

\item 设 \(R\) 是 Dedekind 整环。

（a）证明由「\(I \cong J\) 作为 \(R\)-模当且仅当 \(I \sim J\)」定义的 \(\sim\) 是 \(R\) 的非零分式理想集合上的等价关系。设 \(\mathcal{C}(R)\) 是相应的 \(R\)-模同构类之集，\([I] \in \mathcal{C}(R)\) 表示包含分式理想 \(I\) 的等价类。

（b）证明乘法 \([I][J] = [I \oplus J]\) 给出良定义的二元运算，使 \(\mathcal{C}(R)\) 成为以 \(1 = [R]\) 为单位元的阿贝尔群。

（c）证明（b）中的阿贝尔群 \(\mathcal{C}(R)\) 同构于 \(R\) 的类群。

\item 若 \(R\) 是 Dedekind 整环、\(I\) 是任意非零理想，证明 \(R/I\) 中只有有限多个理想。特别地证明 \(R/I\) 是 Artin 环。

\item 设 \(I\) 是 Dedekind 整环 \(R\) 中的非零分式理想。显式地把 \(I\) 表示为自由 \(R\)-模的直和项，以证明 \(I\) 投射。〔考虑 \(I \oplus I^{-1}\) 并利用命题 21。〕

\item 设 \(I\) 与 \(J\) 是 Dedekind 整环 \(R\) 中的两个非零分式理想，且对某个 \(n \neq 0\) 有 \(I^n = J^n\). 证明 \(I = J\).

\item 设 \(K\) 是代数数域，\(\mathcal{O}_K\) 是 \(K\) 中的整数环。若 \(P\) 是 \(\mathcal{O}_K\) 中的非零素理想，证明对某个素数 \(p \in \mathbb{Z}\) 与某个代数整数 \(\pi \in \mathcal{O}_K\) 有 \(P = (p, \pi)\).

\item 设 \(K = \mathbb{Q}(\sqrt{D})\) 是 \(\mathbb{Q}\) 的二次扩张（\(D\) 是无平方因子整数），\(\mathcal{O}_K\) 是 \(K\) 中的整数环。

（a）证明 \(|\mathcal{O}_K/(p)| \leq p^2\).〔注意作为阿贝尔群 \(\mathcal{O}_K \cong \mathbb{Z}^2\).〕

（b）用推论 16 证明 \((p)\) 在 \(\mathcal{O}_K\) 中的素理想分解有三种可能：

（i）\((p) = P\) 是满足 \(|\mathcal{O}_K/P| = p^2\) 的素理想；

（ii）\((p) = P_1P_2\)，其中 \(P_1, P_2\) 是互不相同的素理想且 \(|\mathcal{O}_K/P_1| = |\mathcal{O}_K/P_2| = p\)；

（iii）\((p) = P^2\)，其中 \(P\) 是满足 \(|\mathcal{O}_K/P| = p\) 的素理想。

（在情形（i）、（ii）、（iii）中分别称素数 \(p\) 在 \(\mathcal{O}_K\) 中\textbf{惯性}、\textbf{分裂}、\textbf{分歧}。分歧素数的集合有限：当 \(D \equiv 1, 2 \pmod 4\) 时是整除 \(D\) 的素数；当 \(D \equiv 3 \pmod 4\) 时是 \(p = 2\) 与整除 \(D\) 的素数。参见第 15.5 节习题 31。）

（c）确定素数 2、3、5、7、11 在 \(K = \mathbb{Q}(\sqrt{-5})\) 的整数环 \(\mathcal{O}_K = \mathbb{Z}[\sqrt{-5}]\) 中的素理想分解。

\item 设 \(\mathcal{O}\) 是 \(\mathbb{Q}\) 的代数闭包 \(\overline{\mathbb{Q}}\) 中的整数环。

（a）证明 \(\mathcal{O}\) 中理想的无限序列 \((2) \subseteq (\sqrt{2}) \subseteq (\sqrt[4]{2}) \subseteq (\sqrt[8]{2}) \subseteq \cdots\) 严格递增，故 \(\mathcal{O}\) 不是诺特环。

（b）证明 \(\mathcal{O}\) 的 Krull 维数为 1.〔利用第 15.3 节定理 26。〕

（c）设 \(K\) 是数域，\(I\) 是 \(\mathcal{O}_K\) 中的任意理想。证明存在 \(K\) 的某个有限扩张 \(L\) 使 \(I\) 扩张到 \(\mathcal{O}_L\) 后成为主理想，即理想 \(I\mathcal{O}_L\) 是主理想（其中 \(L\) 依赖于 \(I\)）——可以使用「\(K\) 的类群是有限群」这一定理。〔参见习题 20。〕

（d）证明 \(\mathcal{O}\) 是 Bézout 整环（参见第 8.1 节）。

\item 设 \(F\) 与 \(K\) 是代数数域，\(\mathbb{Q} \subseteq F \subseteq K\)，整数环分别为 \(\mathcal{O}_F\) 与 \(\mathcal{O}_K\). 由于 \(\mathcal{O}_F \subseteq \mathcal{O}_K\)，环 \(\mathcal{O}_K\) 自然是 \(\mathcal{O}_F\)-模。

（a）证明 \(\mathcal{O}_K\) 是秩为 \(n = [K:F]\) 的无挠 \(\mathcal{O}_F\)-模。〔对 \(\mathbb{Z}\) 计算秩。〕若 \(\mathcal{O}_K\) 在 \(\mathcal{O}_F\) 上自由，则称 \(\mathcal{O}_K\) 在 \(\mathcal{O}_F\) 上有\textbf{相对整基}。

（b）证明若 \(F\) 的类数为 1，则 \(\mathcal{O}_K\) 在 \(\mathcal{O}_F\) 上有相对整基。

若 \(K = \mathbb{Q}(\sqrt{-5}, \sqrt{2})\)，则其整数环由 \(\mathcal{O}_K = \mathbb{Z}+\mathbb{Z}\sqrt{2}+\mathbb{Z}\sqrt{-10}+\mathbb{Z}\omega\) 给出，其中 \(\omega = (\sqrt{-10}+\sqrt{2})/2\).

（c）若 \(F_1 = \mathbb{Q}(\sqrt{2})\)，证明 \(\mathcal{O}_K\) 在 \(\mathcal{O}_{F_1}\) 上有相对整基，并求出显式的基 \(\{\alpha, \beta\}\)：\(\mathcal{O}_K = \mathcal{O}_{F_1}\cdot\alpha+\mathcal{O}_{F_1}\cdot\beta\).

（d）若 \(F_2 = \mathbb{Q}(\sqrt{-5})\)，证明 \(P_3 = (3, 1+\sqrt{-5}) = (3, 5-\sqrt{-5})\) 是 \(\mathcal{O}_{F_2}\) 中非主的素理想，且 \(\mathcal{O}_K = \mathcal{O}_{F_2}\cdot 1+(1/3)P_3\cdot\omega\).〔验证 \(\sqrt{-10} = -(5-\sqrt{-5})\omega/3\).〕由此证明作为 \(\mathcal{O}_{F_2}\)-模，\(\mathcal{O}_K\) 的 Steinitz 类是 \(\mathcal{O}_{F_2}\) 的类群中 \(P_3\) 的非平凡类，故 \(\mathcal{O}_K\) 在 \(\mathcal{O}_{F_2}\) 上没有相对整基。

（e）判断 \(\mathcal{O}_K\) 在 \(K\) 的剩余二次子域 \(F_3 = \mathbb{Q}(\sqrt{-10})\) 的整数环上是否有相对整基。

\item 设 \(C\) 是代数闭域 \(k\) 上的非奇异不可约仿射曲线。证明坐标环 \(k[C]\) 是 Dedekind 整环。
\end{enumerate}
